\documentclass[aspectratio=169,10pt]{beamer}
\usetheme{Madrid}
\usecolortheme{default}
\usepackage[utf8]{inputenc}
\usepackage[T1]{fontenc}
\usepackage{lmodern}
\usepackage{amsmath,amssymb}
\usepackage{enumitem}
\usepackage{xcolor}
\setbeamertemplate{navigation symbols}{}
\setbeamertemplate{footline}[frame number]
\setbeamertemplate{itemize items}[circle]
\setbeamerfont{frametitle}{size=\large,series=\bfseries}
\title{CIA Part 1 -- Final 100 Hard Questions}
\subtitle{Foundations, Ethics, Governance/Risk/Control, Fraud (with IT and AI audit focus)}
\author{Study Deck}
\date{\today}
\begin{document}
\begin{frame}\titlepage\end{frame}
\begin{frame}{Format and blueprint}
\small
\begin{itemize}
\item 100 questions, one per slide; the explanation follows each question (why correct, why the others are wrong).
\item Aligned to the four-domain Part 1 structure: Foundations (35), Ethics and Professionalism (20), Governance, Risk Management and Control (30), Fraud Risks (15).
\item Terminology follows the Global Internal Audit Standards (board, CAE, advisory services, Three Lines Model).
\item About 17 items are technical: ITGC, SOC reports, cloud, resilience, cybersecurity, data analytics, AI governance.
\item All 100 are new and do not repeat the first final set or your original files.
\end{itemize}
\end{frame}

% ===== Q1 =====
\begin{frame}[t]{Q1 \hfill {\scriptsize Foundations of Internal Auditing}}
	\small
	\textbf{Question:} Quality control staff embedded in a production plant, reporting to the plant manager, perform daily checks of finished goods. In the Three Lines Model this activity is:
	\begin{enumerate}[label=\Alph*.,itemsep=1pt,topsep=3pt]
			\item First line.
			\item Second line.
			\item Third line.
			\item Part of the governing body's responsibilities.
	\end{enumerate}
\end{frame}

\begin{frame}[t]{Q1 -- Explanation}
	\footnotesize
	\textbf{Correct: A.} Activities integral to delivering products and managing operational risk under line management belong to the first line. \\[5pt]
\textbf{B wrong:} The second line provides risk and compliance expertise and oversight separate from operations. \\[3pt]
\textbf{C wrong:} The third line is independent internal audit. \\[3pt]
\textbf{D wrong:} The governing body provides oversight, not daily operational checks. \\[3pt]
	\vfill {\tiny Set 2}
\end{frame}

% ===== Q2 =====
\begin{frame}[t]{Q2 \hfill {\scriptsize Foundations of Internal Auditing}}
	\small
	\textbf{Question:} Which is MOST consistent with an agile approach to internal auditing under the Standards?
	\begin{enumerate}[label=\Alph*.,itemsep=1pt,topsep=3pt]
			\item Short iterations with frequent stakeholder communication, while still meeting requirements for supervision, evidence, and documentation.
			\item Dropping documentation to save time.
			\item Eliminating supervisory review for small engagements.
			\item Skipping risk assessment because iterations are short.
	\end{enumerate}
\end{frame}

\begin{frame}[t]{Q2 -- Explanation}
	\footnotesize
	\textbf{Correct: A.} Agile methods can change the delivery approach but not the requirements for quality work. \\[5pt]
\textbf{B wrong:} Documentation requirements still apply. \\[3pt]
\textbf{C wrong:} Supervision remains required. \\[3pt]
\textbf{D wrong:} Risk assessment remains the basis for planning. \\[3pt]
	\vfill {\tiny Set 2}
\end{frame}

% ===== Q3 =====
\begin{frame}[t]{Q3 \hfill {\scriptsize Foundations of Internal Auditing}}
	\small
	\textbf{Question:} The board directs the CAE to remove a high-risk area from the plan to save cost. What should the CAE do?
	\begin{enumerate}[label=\Alph*.,itemsep=1pt,topsep=3pt]
			\item Cover the area anyway without telling the board.
			\item Explain the risk of not covering the area and ensure the board's decision and the accepted risk are documented.
			\item Remove the area without comment.
			\item Resign immediately.
	\end{enumerate}
\end{frame}

\begin{frame}[t]{Q3 -- Explanation}
	\footnotesize
	\textbf{Correct: B.} The CAE informs the board of the consequences; the board's decision and the resulting risk acceptance are recorded. \\[5pt]
\textbf{A wrong:} Acting against the board's decision undermines governance. \\[3pt]
\textbf{C wrong:} Silence leaves the board unaware of the exposure. \\[3pt]
\textbf{D wrong:} Resignation is premature. \\[3pt]
	\vfill {\tiny Set 2}
\end{frame}

% ===== Q4 =====
\begin{frame}[t]{Q4 \hfill {\scriptsize Foundations of Internal Auditing}}
	\small
	\textbf{Question:} Which engagement is an example of an ASSURANCE service?
	\begin{enumerate}[label=\Alph*.,itemsep=1pt,topsep=3pt]
			\item Evaluating the effectiveness of controls over third-party onboarding and giving the board a conclusion.
			\item Facilitating a workshop in which managers identify risks in their own units.
			\item Coaching a project team on control concepts at its request.
			\item Providing observations on a draft policy at management's request, without a conclusion.
	\end{enumerate}
\end{frame}

\begin{frame}[t]{Q4 -- Explanation}
	\footnotesize
	\textbf{Correct: A.} Assurance involves an independent evaluation producing a conclusion for a party other than the process owner. \\[5pt]
\textbf{B wrong:} Facilitation is advisory. \\[3pt]
\textbf{C wrong:} Coaching is advisory. \\[3pt]
\textbf{D wrong:} Observations on a draft policy without a conclusion are advisory. \\[3pt]
	\vfill {\tiny Set 2}
\end{frame}

% ===== Q5 =====
\begin{frame}[t]{Q5 \hfill {\scriptsize Foundations of Internal Auditing}}
	\small
	\textbf{Question:} Management disagrees with the CAE's risk ranking of an area in the internal audit plan. What is the MOST appropriate action?
	\begin{enumerate}[label=\Alph*.,itemsep=1pt,topsep=3pt]
			\item Remove the area from the plan.
			\item Rank it as management prefers.
			\item Ask the external auditor to arbitrate.
			\item Discuss with management; if unresolved, the CAE retains responsibility and presents the risk-based plan to the board.
	\end{enumerate}
\end{frame}

\begin{frame}[t]{Q5 -- Explanation}
	\footnotesize
	\textbf{Correct: D.} The plan is the CAE's responsibility, informed by management input and approved by the board. \\[5pt]
\textbf{A wrong:} That lets management control the plan. \\[3pt]
\textbf{B wrong:} That abandons the CAE's risk assessment. \\[3pt]
\textbf{C wrong:} The external auditor has no authority over the plan. \\[3pt]
	\vfill {\tiny Set 2}
\end{frame}

% ===== Q6 =====
\begin{frame}[t]{Q6 \hfill {\scriptsize Foundations of Internal Auditing}}
	\small
	\textbf{Question:} Who is responsible for developing, implementing, and maintaining the quality assurance and improvement program?
	\begin{enumerate}[label=\Alph*.,itemsep=1pt,topsep=3pt]
			\item The chair of the audit committee.
			\item The external assessor.
			\item The CAE, with oversight by the board.
			\item Each auditor independently, with no coordination.
	\end{enumerate}
\end{frame}

\begin{frame}[t]{Q6 -- Explanation}
	\footnotesize
	\textbf{Correct: C.} The CAE owns the program and reports on it, and the board oversees it. \\[5pt]
\textbf{A wrong:} The board oversees but does not run the program. \\[3pt]
\textbf{B wrong:} The external assessor evaluates at least every five years but does not own the program. \\[3pt]
\textbf{D wrong:} A program requires coordinated management. \\[3pt]
	\vfill {\tiny Set 2}
\end{frame}

% ===== Q7 =====
\begin{frame}[t]{Q7 \hfill {\scriptsize Foundations of Internal Auditing}}
	\small
	\textbf{Question:} Management asks internal audit to review the organization's sustainability reporting. What is the MOST appropriate scope?
	\begin{enumerate}[label=\Alph*.,itemsep=1pt,topsep=3pt]
			\item Give an opinion on the scientific validity of climate targets.
			\item Evaluate governance, risk management, and controls over the collection and reporting of sustainability data.
			\item Prepare the sustainability disclosures.
			\item Set the organization's emissions targets.
	\end{enumerate}
\end{frame}

\begin{frame}[t]{Q7 -- Explanation}
	\footnotesize
	\textbf{Correct: B.} Internal audit assesses processes and controls over reported information without taking management responsibilities. \\[5pt]
\textbf{A wrong:} That requires scientific expertise beyond the question of process and control. \\[3pt]
\textbf{C wrong:} Preparing disclosures is a management responsibility. \\[3pt]
\textbf{D wrong:} Setting targets is a management decision. \\[3pt]
	\vfill {\tiny Set 2}
\end{frame}

% ===== Q8 =====
\begin{frame}[t]{Q8 \hfill {\scriptsize Foundations of Internal Auditing}}
	\small
	\textbf{Question:} Compared with a cyclical audit plan, a risk-based plan:
	\begin{enumerate}[label=\Alph*.,itemsep=1pt,topsep=3pt]
			\item Audits every area at the same interval.
			\item Is set according to management's preferences.
			\item Prioritizes coverage by assessed risk, so higher-risk areas may be audited more often.
			\item Excludes advisory services.
	\end{enumerate}
\end{frame}

\begin{frame}[t]{Q8 -- Explanation}
	\footnotesize
	\textbf{Correct: C.} Risk-based planning allocates resources by risk rather than rotation. \\[5pt]
\textbf{A wrong:} That describes a simple cyclical plan. \\[3pt]
\textbf{B wrong:} Management preference is not a risk assessment. \\[3pt]
\textbf{D wrong:} Advisory services may be included. \\[3pt]
	\vfill {\tiny Set 2}
\end{frame}

% ===== Q9 =====
\begin{frame}[t]{Q9 \hfill {\scriptsize Foundations of Internal Auditing}}
	\small
	\textbf{Question:} When may a report state that work was conducted in conformance with the Standards?
	\begin{enumerate}[label=\Alph*.,itemsep=1pt,topsep=3pt]
			\item Whenever the CAE believes the work was good.
			\item Only when the results of the quality program support the statement.
			\item Always, since the Standards are mandatory.
			\item If the external auditor agrees.
	\end{enumerate}
\end{frame}

\begin{frame}[t]{Q9 -- Explanation}
	\footnotesize
	\textbf{Correct: B.} A conformance claim must be supportable by the results of the quality assurance and improvement program. \\[5pt]
\textbf{A wrong:} Belief alone is not support. \\[3pt]
\textbf{C wrong:} Being mandatory does not mean conformance is assured. \\[3pt]
\textbf{D wrong:} The external auditor does not assess internal audit conformance. \\[3pt]
	\vfill {\tiny Set 2}
\end{frame}

% ===== Q10 =====
\begin{frame}[t]{Q10 \hfill {\scriptsize Foundations of Internal Auditing}}
	\small
	\textbf{Question:} A review of control design for a proposed system, performed while the system is still being designed, is BEST classified as:
	\begin{enumerate}[label=\Alph*.,itemsep=1pt,topsep=3pt]
			\item An advisory service.
			\item An assurance engagement on operating effectiveness.
			\item A compliance audit.
			\item A financial statement audit.
	\end{enumerate}
\end{frame}

\begin{frame}[t]{Q10 -- Explanation}
	\footnotesize
	\textbf{Correct: A.} Pre-implementation design input is advice to the project, with management keeping decision responsibility. \\[5pt]
\textbf{B wrong:} Operating effectiveness cannot be assessed before the system operates. \\[3pt]
\textbf{C wrong:} No compliance requirement is being assessed. \\[3pt]
\textbf{D wrong:} No financial statements are being audited. \\[3pt]
	\vfill {\tiny Set 2}
\end{frame}

% ===== Q11 =====
\begin{frame}[t]{Q11 \hfill {\scriptsize Foundations of Internal Auditing}}
	\small
	\textbf{Question:} A multinational has regional internal audit teams. Which approach BEST supports consistent quality?
	\begin{enumerate}[label=\Alph*.,itemsep=1pt,topsep=3pt]
			\item Each region sets its own methodology.
			\item Regional CFOs supervise the local teams.
			\item Quality is assessed only at group level every five years.
			\item The group CAE sets common methodology and supervises and reviews work across regions.
	\end{enumerate}
\end{frame}

\begin{frame}[t]{Q11 -- Explanation}
	\footnotesize
	\textbf{Correct: D.} Consistent methodology and supervision are CAE responsibilities across the function. \\[5pt]
\textbf{A wrong:} Divergent methodologies create inconsistent quality. \\[3pt]
\textbf{B wrong:} CFO supervision would impair independence. \\[3pt]
\textbf{C wrong:} Periodic assessment alone is not ongoing quality management. \\[3pt]
	\vfill {\tiny Set 2}
\end{frame}

% ===== Q12 =====
\begin{frame}[t]{Q12 \hfill {\scriptsize Foundations of Internal Auditing}}
	\small
	\textbf{Question:} A new internal audit function is being established. What should the CAE do FIRST?
	\begin{enumerate}[label=\Alph*.,itemsep=1pt,topsep=3pt]
			\item Draft a charter defining mandate and reporting lines and obtain board approval.
			\item Hire staff without defining the mandate.
			\item Prepare detailed engagement programs.
			\item Obtain the external auditor's approval of the audit universe.
	\end{enumerate}
\end{frame}

\begin{frame}[t]{Q12 -- Explanation}
	\footnotesize
	\textbf{Correct: A.} The charter establishes authority, role, and reporting relationships on which the rest depends. \\[5pt]
\textbf{B wrong:} Without a mandate, resourcing lacks direction. \\[3pt]
\textbf{C wrong:} Engagement programs depend on a plan and mandate. \\[3pt]
\textbf{D wrong:} The external auditor has no approval role. \\[3pt]
	\vfill {\tiny Set 2}
\end{frame}

% ===== Q13 =====
\begin{frame}[t]{Q13 \hfill {\scriptsize Foundations of Internal Auditing}}
	\small
	\textbf{Question:} Which statement BEST reflects the mission of internal audit under the Global Internal Audit Standards?
	\begin{enumerate}[label=\Alph*.,itemsep=1pt,topsep=3pt]
			\item Help management meet its targets by taking responsibility for risk decisions and control design.
			\item Provide the external auditor with the evidence needed to support the financial statement opinion.
			\item Enhance and protect organizational value by providing risk-based and objective assurance, advice, and insight.
			\item Enforce compliance with policies by investigating and disciplining violations.
	\end{enumerate}
\end{frame}

\begin{frame}[t]{Q13 -- Explanation}
	\footnotesize
	\textbf{Correct: C.} The mission is to enhance and protect organizational value through risk-based, objective assurance, advice, and insight. \\[5pt]
\textbf{A wrong:} Internal audit must not take management's decision-making responsibilities. \\[3pt]
\textbf{B wrong:} Supporting the external audit is not the mission; internal audit serves the board and management. \\[3pt]
\textbf{D wrong:} Enforcement and discipline are management responsibilities. \\[3pt]
	\vfill {\tiny Set 2}
\end{frame}

% ===== Q14 =====
\begin{frame}[t]{Q14 \hfill {\scriptsize Foundations of Internal Auditing}}
	\small
	\textbf{Question:} Management asks internal audit to help set the company's new strategic objectives. What is the MOST appropriate role?
	\begin{enumerate}[label=\Alph*.,itemsep=1pt,topsep=3pt]
			\item Define the objectives, since internal audit knows the business best.
			\item Approve the objectives before the board reviews them.
			\item Refuse any involvement in strategy-related matters.
			\item Provide assurance and advice on governance and risk considerations in the strategy-setting process, without setting the objectives.
	\end{enumerate}
\end{frame}

\begin{frame}[t]{Q14 -- Explanation}
	\footnotesize
	\textbf{Correct: D.} Internal audit can provide insight and advice but must not assume management's decision-making role. \\[5pt]
\textbf{A wrong:} Setting objectives is a management responsibility. \\[3pt]
\textbf{B wrong:} Approval would be a management or governance decision. \\[3pt]
\textbf{C wrong:} Advisory involvement is permitted when objectivity is preserved. \\[3pt]
	\vfill {\tiny Set 2}
\end{frame}

% ===== Q15 =====
\begin{frame}[t]{Q15 \hfill {\scriptsize Foundations of Internal Auditing}}
	\small
	\textbf{Question:} A department head refuses access to records, citing sensitivity, and the CAE cannot resolve this with him. What should the CAE do?
	\begin{enumerate}[label=\Alph*.,itemsep=1pt,topsep=3pt]
			\item Close the engagement without comment.
			\item Substitute interview evidence and not mention the restriction.
			\item Ask the external auditor to obtain the records.
			\item Escalate to senior management and, if still unresolved, to the board, because it is a restriction on access.
	\end{enumerate}
\end{frame}

\begin{frame}[t]{Q15 -- Explanation}
	\footnotesize
	\textbf{Correct: D.} Restrictions on access affect the mandate and must be escalated until resolved or the board is informed. \\[5pt]
\textbf{A wrong:} Closing silently hides a limitation. \\[3pt]
\textbf{B wrong:} Not disclosing the limitation misleads readers of the results. \\[3pt]
\textbf{C wrong:} The external auditor cannot obtain the records on internal audit's behalf. \\[3pt]
	\vfill {\tiny Set 2}
\end{frame}

% ===== Q16 =====
\begin{frame}[t]{Q16 \hfill {\scriptsize Foundations of Internal Auditing}}
	\small
	\textbf{Question:} How should the CAE plan for emerging risks that arise between planning cycles?
	\begin{enumerate}[label=\Alph*.,itemsep=1pt,topsep=3pt]
			\item Fix the plan for the year to protect independence.
			\item Rely on the external auditors to identify emerging risks.
			\item Add emerging risks only at the next annual cycle.
			\item Include flexible capacity and reassess risks periodically so the plan can be adjusted.
	\end{enumerate}
\end{frame}

\begin{frame}[t]{Q16 -- Explanation}
	\footnotesize
	\textbf{Correct: D.} A risk-based plan must stay current, with capacity to respond to change. \\[5pt]
\textbf{A wrong:} A rigid plan may miss significant risks. \\[3pt]
\textbf{B wrong:} External auditors have a different focus. \\[3pt]
\textbf{C wrong:} Waiting for the next cycle leaves new risks uncovered. \\[3pt]
	\vfill {\tiny Set 2}
\end{frame}

% ===== Q17 =====
\begin{frame}[t]{Q17 \hfill {\scriptsize Foundations of Internal Auditing}}
	\small
	\textbf{Question:} Which authority must the internal audit mandate provide so the function can carry out its work?
	\begin{enumerate}[label=\Alph*.,itemsep=1pt,topsep=3pt]
			\item Authority to direct employees to change their processes.
			\item Unrestricted access to the records, personnel, and physical properties relevant to engagements.
			\item Authority to approve management's budgets.
			\item Authority to approve risk responses.
	\end{enumerate}
\end{frame}

\begin{frame}[t]{Q17 -- Explanation}
	\footnotesize
	\textbf{Correct: B.} Access to information, people, and property is essential for independent evaluation. \\[5pt]
\textbf{A wrong:} Directing process changes is a management role. \\[3pt]
\textbf{C wrong:} Approving budgets is a management role. \\[3pt]
\textbf{D wrong:} Approving risk responses is a management responsibility. \\[3pt]
	\vfill {\tiny Set 2}
\end{frame}

% ===== Q18 =====
\begin{frame}[t]{Q18 \hfill {\scriptsize Foundations of Internal Auditing}}
	\small
	\textbf{Question:} What is the main benefit of including root-cause analysis in findings?
	\begin{enumerate}[label=\Alph*.,itemsep=1pt,topsep=3pt]
			\item It identifies the underlying reason for the condition so recommendations address the cause, not symptoms.
			\item It assigns blame to individuals.
			\item It reduces the need for criteria.
			\item It guarantees that management agrees.
	\end{enumerate}
\end{frame}

\begin{frame}[t]{Q18 -- Explanation}
	\footnotesize
	\textbf{Correct: A.} Understanding the cause helps produce more effective, lasting solutions. \\[5pt]
\textbf{B wrong:} Root-cause analysis focuses on process and system causes, not blame. \\[3pt]
\textbf{C wrong:} Criteria remain necessary. \\[3pt]
\textbf{D wrong:} Agreement is not guaranteed. \\[3pt]
	\vfill {\tiny Set 2}
\end{frame}

% ===== Q19 =====
\begin{frame}[t]{Q19 \hfill {\scriptsize Foundations of Internal Auditing}}
	\small
	\textbf{Question:} The CAE proposes shifting the function from mostly compliance work to more advisory and strategic support. What is required?
	\begin{enumerate}[label=\Alph*.,itemsep=1pt,topsep=3pt]
			\item Obtain the CEO's verbal agreement.
			\item Update the charter and plan to reflect the change, with board approval and a review of resources and competencies.
			\item No action, since the Standards cover both types.
			\item Obtain the external auditor's consent.
	\end{enumerate}
\end{frame}

\begin{frame}[t]{Q19 -- Explanation}
	\footnotesize
	\textbf{Correct: B.} A change in mandate or focus affects the charter, plan, and resources, which the board approves. \\[5pt]
\textbf{A wrong:} A verbal agreement from the CEO does not change the mandate. \\[3pt]
\textbf{C wrong:} The mandate and plan still need to be updated. \\[3pt]
\textbf{D wrong:} The external auditor has no role. \\[3pt]
	\vfill {\tiny Set 2}
\end{frame}

% ===== Q20 =====
\begin{frame}[t]{Q20 \hfill {\scriptsize Foundations of Internal Auditing}}
	\small
	\textbf{Question:} A report omits information that readers need to understand its conclusions. Which communication quality is MOST lacking?
	\begin{enumerate}[label=\Alph*.,itemsep=1pt,topsep=3pt]
			\item Completeness.
			\item Conciseness.
			\item Timeliness.
			\item Objectivity.
	\end{enumerate}
\end{frame}

\begin{frame}[t]{Q20 -- Explanation}
	\footnotesize
	\textbf{Correct: A.} Complete communications include all relevant information needed by the intended users. \\[5pt]
\textbf{B wrong:} Conciseness means avoiding unnecessary detail, not omitting necessary information. \\[3pt]
\textbf{C wrong:} Timeliness concerns when the report is delivered. \\[3pt]
\textbf{D wrong:} Objectivity concerns fairness and lack of bias. \\[3pt]
	\vfill {\tiny Set 2}
\end{frame}

% ===== Q21 =====
\begin{frame}[t]{Q21 \hfill {\scriptsize Foundations of Internal Auditing}}
	\small
	\textbf{Question:} Which information should the CAE report to the board periodically?
	\begin{enumerate}[label=\Alph*.,itemsep=1pt,topsep=3pt]
			\item Only the list of findings closed.
			\item Only issues that management agrees are significant.
			\item Performance against the plan, conformance with the Standards, significant risk and control issues, and action plan status.
			\item Only budget utilization.
	\end{enumerate}
\end{frame}

\begin{frame}[t]{Q21 -- Explanation}
	\footnotesize
	\textbf{Correct: C.} The board needs a complete view of the function's performance, quality, and significant matters. \\[5pt]
\textbf{A wrong:} That is too narrow. \\[3pt]
\textbf{B wrong:} Management agreement should not filter board reporting. \\[3pt]
\textbf{D wrong:} Budget alone does not support oversight. \\[3pt]
	\vfill {\tiny Set 2}
\end{frame}

% ===== Q22 =====
\begin{frame}[t]{Q22 \hfill {\scriptsize Foundations of Internal Auditing}}
	\small
	\textbf{Question:} Which statement BEST distinguishes internal audit from external audit?
	\begin{enumerate}[label=\Alph*.,itemsep=1pt,topsep=3pt]
			\item Internal audit expresses the opinion on financial statements for shareholders.
			\item Internal audit serves the board and management with a broad focus on governance, risk, and control; external audit primarily supports an opinion on financial statements for external users.
			\item External audit reports to the CAE.
			\item They have identical scope, differing only in staff.
	\end{enumerate}
\end{frame}

\begin{frame}[t]{Q22 -- Explanation}
	\footnotesize
	\textbf{Correct: B.} The two functions differ in audience, scope, and objectives, though they coordinate. \\[5pt]
\textbf{A wrong:} External auditors express the financial statement opinion. \\[3pt]
\textbf{C wrong:} External auditors do not report to the CAE. \\[3pt]
\textbf{D wrong:} Their scope and purpose differ. \\[3pt]
	\vfill {\tiny Set 2}
\end{frame}

% ===== Q23 =====
\begin{frame}[t]{Q23 \hfill {\scriptsize Foundations of Internal Auditing}}
	\small
	\textbf{Question:} An auditor leads a workshop in which managers identify and rate risks in their own units, and management keeps ownership of the results. This is BEST classified as:
	\begin{enumerate}[label=\Alph*.,itemsep=1pt,topsep=3pt]
			\item An advisory service (facilitation).
			\item An assurance engagement.
			\item An investigation.
			\item An external audit procedure.
	\end{enumerate}
\end{frame}

\begin{frame}[t]{Q23 -- Explanation}
	\footnotesize
	\textbf{Correct: A.} Facilitation helps management while management retains ownership of the content and decisions. \\[5pt]
\textbf{B wrong:} No independent conclusion is provided to a third party. \\[3pt]
\textbf{C wrong:} There is no allegation being examined. \\[3pt]
\textbf{D wrong:} It is not part of an external audit. \\[3pt]
	\vfill {\tiny Set 2}
\end{frame}

% ===== Q24 =====
\begin{frame}[t]{Q24 \hfill {\scriptsize Foundations of Internal Auditing}}
	\small
	\textbf{Question:} Which statement correctly distinguishes continuous monitoring from continuous auditing?
	\begin{enumerate}[label=\Alph*.,itemsep=1pt,topsep=3pt]
			\item Internal audit performs continuous monitoring and management performs continuous auditing.
			\item Management performs continuous monitoring of controls and risks; internal audit performs continuous auditing to independently assess them.
			\item Both are performed only by external auditors.
			\item Continuous auditing replaces the need for any management monitoring.
	\end{enumerate}
\end{frame}

\begin{frame}[t]{Q24 -- Explanation}
	\footnotesize
	\textbf{Correct: B.} Monitoring is a management responsibility; continuous auditing is the independent assurance counterpart. \\[5pt]
\textbf{A wrong:} That reverses the roles. \\[3pt]
\textbf{C wrong:} Neither is limited to external auditors. \\[3pt]
\textbf{D wrong:} Management remains responsible for monitoring its controls. \\[3pt]
	\vfill {\tiny Set 2}
\end{frame}

% ===== Q25 =====
\begin{frame}[t]{Q25 \hfill {\scriptsize Foundations of Internal Auditing}}
	\small
	\textbf{Question:} Which of the following is a responsibility of the board, as opposed to the CAE or management, in relation to internal audit?
	\begin{enumerate}[label=\Alph*.,itemsep=1pt,topsep=3pt]
			\item Designing the audit program for each engagement.
			\item Approving the risk-based internal audit plan and the internal audit resource plan.
			\item Selecting the sampling method for testing.
			\item Deciding the wording of engagement observations.
	\end{enumerate}
\end{frame}

\begin{frame}[t]{Q25 -- Explanation}
	\footnotesize
	\textbf{Correct: B.} The board approves the plan and resource plan, which supports independence and oversight of the function. \\[5pt]
\textbf{A wrong:} Engagement programs are designed by internal audit. \\[3pt]
\textbf{C wrong:} Testing methods are an engagement-level decision for internal audit. \\[3pt]
\textbf{D wrong:} Observation wording is determined by internal audit, not the board. \\[3pt]
	\vfill {\tiny Set 2}
\end{frame}

% ===== Q26 =====
\begin{frame}[t]{Q26 \hfill {\scriptsize Foundations of Internal Auditing}}
	\small
	\textbf{Question:} A manager from the sales department joins internal audit for a six-month rotation. Which first assignment is MOST appropriate?
	\begin{enumerate}[label=\Alph*.,itemsep=1pt,topsep=3pt]
			\item An assurance engagement on sales commissions.
			\item An assurance engagement on the sales order process.
			\item An audit of a function unrelated to sales, with supervision.
			\item A review of his former team's controls.
	\end{enumerate}
\end{frame}

\begin{frame}[t]{Q26 -- Explanation}
	\footnotesize
	\textbf{Correct: C.} Objectivity is presumed impaired for assurance on areas the auditor was responsible for within the prior 12 months. \\[5pt]
\textbf{A wrong:} He was recently responsible for sales activities. \\[3pt]
\textbf{B wrong:} The same presumption applies. \\[3pt]
\textbf{D wrong:} That is the area he just left. \\[3pt]
	\vfill {\tiny Set 2}
\end{frame}

% ===== Q27 =====
\begin{frame}[t]{Q27 \hfill {\scriptsize Foundations of Internal Auditing}}
	\small
	\textbf{Question:} Which person would be disqualified from serving on an external quality assessment team?
	\begin{enumerate}[label=\Alph*.,itemsep=1pt,topsep=3pt]
			\item A consultant who designed the function's audit methodology last year.
			\item An experienced CIA from another organization with no ties to the function.
			\item A professional trained in quality assessment methodology.
			\item A CAE from an unrelated company who is competent in the Standards.
	\end{enumerate}
\end{frame}

\begin{frame}[t]{Q27 -- Explanation}
	\footnotesize
	\textbf{Correct: A.} Assessors must be independent and free of conflicts; reviewing one's own design is a self-review conflict. \\[5pt]
\textbf{B wrong:} Independent and qualified. \\[3pt]
\textbf{C wrong:} Qualified with no conflict indicated. \\[3pt]
\textbf{D wrong:} Independent and competent. \\[3pt]
	\vfill {\tiny Set 2}
\end{frame}

% ===== Q28 =====
\begin{frame}[t]{Q28 \hfill {\scriptsize Foundations of Internal Auditing}}
	\small
	\textbf{Question:} A subsidiary's CFO proposes that the local internal audit team report to him. What is the MOST appropriate response from the group CAE?
	\begin{enumerate}[label=\Alph*.,itemsep=1pt,topsep=3pt]
			\item Accept, since the CFO knows the subsidiary best.
			\item Disband the local team.
			\item Let the subsidiary board decide without group involvement.
			\item Maintain functional reporting of the local team to the group CAE and board, with administrative arrangements agreed locally.
	\end{enumerate}
\end{frame}

\begin{frame}[t]{Q28 -- Explanation}
	\footnotesize
	\textbf{Correct: D.} Functional reporting to the CAE and board protects independence; administrative support may be local. \\[5pt]
\textbf{A wrong:} Reporting to the CFO would impair independence. \\[3pt]
\textbf{B wrong:} Disbanding is unnecessary. \\[3pt]
\textbf{C wrong:} The group CAE's role over the function must be maintained. \\[3pt]
	\vfill {\tiny Set 2}
\end{frame}

% ===== Q29 =====
\begin{frame}[t]{Q29 \hfill {\scriptsize Foundations of Internal Auditing}}
	\small
	\textbf{Question:} Management has not implemented an agreed action plan for a high-risk finding by the due date. What should the CAE do FIRST?
	\begin{enumerate}[label=\Alph*.,itemsep=1pt,topsep=3pt]
			\item Close the finding because the due date has passed.
			\item Report management to the regulator.
			\item Understand the reason for the delay and discuss status with management, escalating according to policy if the risk remains unaddressed.
			\item Issue a new finding on management's competence.
	\end{enumerate}
\end{frame}

\begin{frame}[t]{Q29 -- Explanation}
	\footnotesize
	\textbf{Correct: C.} Follow-up starts with understanding the cause, then escalation if the risk persists. \\[5pt]
\textbf{A wrong:} Closing leaves the risk unresolved. \\[3pt]
\textbf{B wrong:} Going to a regulator skips internal escalation channels. \\[3pt]
\textbf{D wrong:} That is premature and not evidence-based. \\[3pt]
	\vfill {\tiny Set 2}
\end{frame}

% ===== Q30 =====
\begin{frame}[t]{Q30 \hfill {\scriptsize Foundations of Internal Auditing}}
	\small
	\textbf{Question:} Who is responsible for ensuring that engagement records are retained in line with organizational policy and legal requirements?
	\begin{enumerate}[label=\Alph*.,itemsep=1pt,topsep=3pt]
			\item The records management department of the auditee.
			\item The CAE.
			\item The external auditor.
			\item The audit committee chair.
	\end{enumerate}
\end{frame}

\begin{frame}[t]{Q30 -- Explanation}
	\footnotesize
	\textbf{Correct: B.} The CAE is responsible for the control of engagement records and their retention. \\[5pt]
\textbf{A wrong:} The auditee does not control audit records. \\[3pt]
\textbf{C wrong:} The external auditor has no role. \\[3pt]
\textbf{D wrong:} The board oversees but does not manage records. \\[3pt]
	\vfill {\tiny Set 2}
\end{frame}

% ===== Q31 =====
\begin{frame}[t]{Q31 \hfill {\scriptsize Foundations of Internal Auditing}}
	\small
	\textbf{Question:} Which is a responsibility of senior management regarding internal audit?
	\begin{enumerate}[label=\Alph*.,itemsep=1pt,topsep=3pt]
			\item Deciding the scope and conclusions of each engagement.
			\item Supporting the function by ensuring access to people and records and cooperation needed to carry out its mandate.
			\item Approving final reports before they go to the board.
			\item Evaluating the CAE's performance without board input.
	\end{enumerate}
\end{frame}

\begin{frame}[t]{Q31 -- Explanation}
	\footnotesize
	\textbf{Correct: B.} Senior management supports the function and its mandate; it does not direct engagement scope or conclusions. \\[5pt]
\textbf{A wrong:} Management directing scope and conclusions would impair independence. \\[3pt]
\textbf{C wrong:} Pre-approval of reports by management would filter board communication. \\[3pt]
\textbf{D wrong:} Performance evaluation of the CAE is a board matter, with management input. \\[3pt]
	\vfill {\tiny Set 2}
\end{frame}

% ===== Q32 =====
\begin{frame}[t]{Q32 \hfill {\scriptsize Foundations of Internal Auditing}}
	\small
	\textbf{Question:} Which of the following is NOT a characteristic of the third line in the Three Lines Model?
	\begin{enumerate}[label=\Alph*.,itemsep=1pt,topsep=3pt]
			\item Providing independent and objective assurance.
			\item Owning and managing risks as part of day-to-day operations.
			\item Reporting to the governing body on its conclusions.
			\item Remaining independent from management's responsibilities.
	\end{enumerate}
\end{frame}

\begin{frame}[t]{Q32 -- Explanation}
	\footnotesize
	\textbf{Correct: B.} Owning and managing risks day to day is a first-line responsibility. \\[5pt]
\textbf{A wrong:} This is a third-line characteristic. \\[3pt]
\textbf{C wrong:} This is a third-line characteristic. \\[3pt]
\textbf{D wrong:} This is a third-line characteristic. \\[3pt]
	\vfill {\tiny Set 2}
\end{frame}

% ===== Q33 =====
\begin{frame}[t]{Q33 \hfill {\scriptsize Foundations of Internal Auditing}}
	\small
	\textbf{Question:} A regulator asks for copies of internal audit reports. What should the CAE do?
	\begin{enumerate}[label=\Alph*.,itemsep=1pt,topsep=3pt]
			\item Follow the established protocol: consult senior management, legal counsel, and the board as appropriate before release, in line with law and policy.
			\item Refuse because reports are confidential.
			\item Release all reports immediately.
			\item Let the audited units decide what to release.
	\end{enumerate}
\end{frame}

\begin{frame}[t]{Q33 -- Explanation}
	\footnotesize
	\textbf{Correct: A.} Disclosure to outside parties should be controlled and consistent with law and policy. \\[5pt]
\textbf{B wrong:} Refusing may violate legal or regulatory requirements. \\[3pt]
\textbf{C wrong:} Unreviewed release may exceed what is required. \\[3pt]
\textbf{D wrong:} Audited units have a conflict of interest. \\[3pt]
	\vfill {\tiny Set 2}
\end{frame}

% ===== Q34 =====
\begin{frame}[t]{Q34 \hfill {\scriptsize Foundations of Internal Auditing}}
	\small
	\textbf{Question:} An engagement was completed with a restriction on information the auditors could obtain. What must the engagement communication include?
	\begin{enumerate}[label=\Alph*.,itemsep=1pt,topsep=3pt]
			\item Nothing, if the auditors believe the remaining evidence is sufficient.
			\item The limitation, but only in the workpapers.
			\item A statement that the scope was fully achieved.
			\item The limitation and its possible effect on the conclusions.
	\end{enumerate}
\end{frame}

\begin{frame}[t]{Q34 -- Explanation}
	\footnotesize
	\textbf{Correct: D.} Transparency about limitations helps users interpret the conclusions. \\[5pt]
\textbf{A wrong:} Board and management cannot assess reliance without knowing about the limitation. \\[3pt]
\textbf{B wrong:} Workpapers are not read by report users. \\[3pt]
\textbf{C wrong:} That would be inaccurate. \\[3pt]
	\vfill {\tiny Set 2}
\end{frame}

% ===== Q35 =====
\begin{frame}[t]{Q35 \hfill {\scriptsize Foundations of Internal Auditing}}
	\small
	\textbf{Question:} Which is the BEST evidence that internal audit adds value for the board?
	\begin{enumerate}[label=\Alph*.,itemsep=1pt,topsep=3pt]
			\item The number of reports issued.
			\item Audit budget as a percentage of revenue.
			\item Recommendations and insights that improve risk management and control outcomes, supported by board and stakeholder feedback.
			\item The number of auditors holding certifications.
	\end{enumerate}
\end{frame}

\begin{frame}[t]{Q35 -- Explanation}
	\footnotesize
	\textbf{Correct: C.} Value is demonstrated by outcomes and stakeholder perception, not by activity or input counts. \\[5pt]
\textbf{A wrong:} Output counts do not show impact. \\[3pt]
\textbf{B wrong:} Cost ratios say nothing about value delivered. \\[3pt]
\textbf{D wrong:} Certifications are an input, not an outcome. \\[3pt]
	\vfill {\tiny Set 2}
\end{frame}

% ===== Q36 =====
\begin{frame}[t]{Q36 \hfill {\scriptsize Ethics and Professionalism}}
	\small
	\textbf{Question:} An auditor reports truthfully on a fraud-related finding despite pressure from senior management to stay silent. Which element of the Code of Ethics is MOST directly demonstrated?
	\begin{enumerate}[label=\Alph*.,itemsep=1pt,topsep=3pt]
			\item Competency.
			\item Confidentiality.
			\item Integrity, including honesty and professional courage.
			\item Cost efficiency.
	\end{enumerate}
\end{frame}

\begin{frame}[t]{Q36 -- Explanation}
	\footnotesize
	\textbf{Correct: C.} Integrity includes honesty and the courage to report even under pressure. \\[5pt]
\textbf{A wrong:} Competency concerns skills. \\[3pt]
\textbf{B wrong:} Confidentiality concerns protection of information. \\[3pt]
\textbf{D wrong:} Cost efficiency is not an ethical principle in the Code. \\[3pt]
	\vfill {\tiny Set 2}
\end{frame}

% ===== Q37 =====
\begin{frame}[t]{Q37 \hfill {\scriptsize Ethics and Professionalism}}
	\small
	\textbf{Question:} The CAE believes management has committed an illegal act and the board has not acted. What is the MOST appropriate next step?
	\begin{enumerate}[label=\Alph*.,itemsep=1pt,topsep=3pt]
			\item Report immediately to the media.
			\item Resign without further action.
			\item Do nothing since the board has decided.
			\item Consult legal counsel on legal and professional obligations and escalate within the governance structure before considering other action.
	\end{enumerate}
\end{frame}

\begin{frame}[t]{Q37 -- Explanation}
	\footnotesize
	\textbf{Correct: D.} The CAE considers legal and professional obligations and pursues escalation through appropriate channels. \\[5pt]
\textbf{A wrong:} Media disclosure breaches confidentiality and is premature. \\[3pt]
\textbf{B wrong:} Resignation does not address the matter. \\[3pt]
\textbf{C wrong:} Inaction ignores a serious legal concern. \\[3pt]
	\vfill {\tiny Set 2}
\end{frame}

% ===== Q38 =====
\begin{frame}[t]{Q38 \hfill {\scriptsize Ethics and Professionalism}}
	\small
	\textbf{Question:} An auditor plans to analyze employee emails with an analytics tool for an engagement. What should happen FIRST?
	\begin{enumerate}[label=\Alph*.,itemsep=1pt,topsep=3pt]
			\item Run the analysis first and seek approval if the results are significant.
			\item Share the dataset with the vendor for faster processing.
			\item Confirm the legal basis, policies, and approvals for using personal data, and limit use to what the engagement requires.
			\item Skip approvals since internal audit has unrestricted access.
	\end{enumerate}
\end{frame}

\begin{frame}[t]{Q38 -- Explanation}
	\footnotesize
	\textbf{Correct: C.} Access rights do not override privacy law and policy; use of personal data must be authorized and proportionate. \\[5pt]
\textbf{A wrong:} Approval should come before processing. \\[3pt]
\textbf{B wrong:} Sharing data externally increases risk. \\[3pt]
\textbf{D wrong:} Unrestricted access to records does not remove legal and policy constraints. \\[3pt]
	\vfill {\tiny Set 2}
\end{frame}

% ===== Q39 =====
\begin{frame}[t]{Q39 \hfill {\scriptsize Ethics and Professionalism}}
	\small
	\textbf{Question:} The audit team uses an external IT specialist. Which statement is MOST accurate regarding due professional care?
	\begin{enumerate}[label=\Alph*.,itemsep=1pt,topsep=3pt]
			\item The CAE must evaluate the specialist's competence and objectivity and remains responsible for the conclusions.
			\item The specialist's conclusions can be adopted without review.
			\item Responsibility transfers to the specialist.
			\item Specialists need not be independent.
	\end{enumerate}
\end{frame}

\begin{frame}[t]{Q39 -- Explanation}
	\footnotesize
	\textbf{Correct: A.} Using a specialist does not transfer accountability for internal audit's conclusions. \\[5pt]
\textbf{B wrong:} Unreviewed reliance is not due care. \\[3pt]
\textbf{C wrong:} Responsibility stays with internal audit. \\[3pt]
\textbf{D wrong:} Objectivity of the specialist must be assessed. \\[3pt]
	\vfill {\tiny Set 2}
\end{frame}

% ===== Q40 =====
\begin{frame}[t]{Q40 \hfill {\scriptsize Ethics and Professionalism}}
	\small
	\textbf{Question:} An auditor had a personal dispute with the manager of the area she is auditing. What is the MOST appropriate action?
	\begin{enumerate}[label=\Alph*.,itemsep=1pt,topsep=3pt]
			\item Continue; being professional is enough.
			\item Continue but avoid the manager entirely.
			\item Resign from the function.
			\item Disclose it to the CAE, who determines whether to reassign or add safeguards.
	\end{enumerate}
\end{frame}

\begin{frame}[t]{Q40 -- Explanation}
	\footnotesize
	\textbf{Correct: D.} Personal bias can impair objectivity and should be disclosed so the CAE can respond. \\[5pt]
\textbf{A wrong:} Self-assessment of bias is not sufficient. \\[3pt]
\textbf{B wrong:} Avoiding the manager limits the work and doesn't remove bias. \\[3pt]
\textbf{C wrong:} Resignation is disproportionate. \\[3pt]
	\vfill {\tiny Set 2}
\end{frame}

% ===== Q41 =====
\begin{frame}[t]{Q41 \hfill {\scriptsize Ethics and Professionalism}}
	\small
	\textbf{Question:} An auditor posts an anonymized description of an ongoing audit finding on social media. Which statement is MOST accurate?
	\begin{enumerate}[label=\Alph*.,itemsep=1pt,topsep=3pt]
			\item It is acceptable because names were removed.
			\item It likely breaches confidentiality, because engagement information should not be disclosed without authority even if the entity is not named.
			\item It is acceptable if it contains no financial figures.
			\item It is acceptable if posted after working hours.
	\end{enumerate}
\end{frame}

\begin{frame}[t]{Q41 -- Explanation}
	\footnotesize
	\textbf{Correct: B.} Auditors protect information obtained in the course of work; anonymization does not authorize disclosure. \\[5pt]
\textbf{A wrong:} Context may still identify the organization. \\[3pt]
\textbf{C wrong:} Non-financial details can also be sensitive. \\[3pt]
\textbf{D wrong:} Timing does not change the obligation. \\[3pt]
	\vfill {\tiny Set 2}
\end{frame}

% ===== Q42 =====
\begin{frame}[t]{Q42 \hfill {\scriptsize Ethics and Professionalism}}
	\small
	\textbf{Question:} Who must conform to the Code of Ethics?
	\begin{enumerate}[label=\Alph*.,itemsep=1pt,topsep=3pt]
			\item Only certified internal auditors.
			\item Only the CAE.
			\item Only auditors performing assurance.
			\item All internal auditors, including the CAE and staff who are not certified.
	\end{enumerate}
\end{frame}

\begin{frame}[t]{Q42 -- Explanation}
	\footnotesize
	\textbf{Correct: D.} The Code applies to all those who provide internal audit services. \\[5pt]
\textbf{A wrong:} Conformance is not limited to certified staff. \\[3pt]
\textbf{B wrong:} It applies to all staff. \\[3pt]
\textbf{C wrong:} It applies to advisory as well. \\[3pt]
	\vfill {\tiny Set 2}
\end{frame}

% ===== Q43 =====
\begin{frame}[t]{Q43 \hfill {\scriptsize Ethics and Professionalism}}
	\small
	\textbf{Question:} Out of curiosity, an auditor shares draft findings with a manager from another department who has no need to know. Which principle is MOST directly violated?
	\begin{enumerate}[label=\Alph*.,itemsep=1pt,topsep=3pt]
			\item Confidentiality.
			\item Competency.
			\item Objectivity.
			\item Sampling methodology.
	\end{enumerate}
\end{frame}

\begin{frame}[t]{Q43 -- Explanation}
	\footnotesize
	\textbf{Correct: A.} Engagement information must be shared only with authorized parties. \\[5pt]
\textbf{B wrong:} Competency concerns skills and knowledge. \\[3pt]
\textbf{C wrong:} Objectivity concerns bias and impairments. \\[3pt]
\textbf{D wrong:} Sampling is a technique and not a principle. \\[3pt]
	\vfill {\tiny Set 2}
\end{frame}

% ===== Q44 =====
\begin{frame}[t]{Q44 \hfill {\scriptsize Ethics and Professionalism}}
	\small
	\textbf{Question:} The auditee proposes the audit program and asks that internal audit follow it. What is the MOST appropriate response?
	\begin{enumerate}[label=\Alph*.,itemsep=1pt,topsep=3pt]
			\item Adopt it to maintain a cooperative relationship.
			\item Adopt it but reduce the testing.
			\item Consider management's input but retain responsibility for scope and procedures based on risk and engagement objectives.
			\item Reject all input from management.
	\end{enumerate}
\end{frame}

\begin{frame}[t]{Q44 -- Explanation}
	\footnotesize
	\textbf{Correct: C.} Internal audit determines scope and procedures independently while taking relevant input. \\[5pt]
\textbf{A wrong:} Adopting auditee-designed procedures risks independence. \\[3pt]
\textbf{B wrong:} Reducing testing weakens the work. \\[3pt]
\textbf{D wrong:} Management input can be valuable. \\[3pt]
	\vfill {\tiny Set 2}
\end{frame}

% ===== Q45 =====
\begin{frame}[t]{Q45 \hfill {\scriptsize Ethics and Professionalism}}
	\small
	\textbf{Question:} The CEO asks the CAE to issue a positive conclusion on a project by Friday, but the audit work will not be finished. What should the CAE do?
	\begin{enumerate}[label=\Alph*.,itemsep=1pt,topsep=3pt]
			\item Issue the conclusion now and finish the work later.
			\item Issue it with boilerplate wording only.
			\item Not issue an unsupported conclusion; limit the communication to work completed and explain the scope limitation.
			\item Ask the auditor to sign without review.
	\end{enumerate}
\end{frame}

\begin{frame}[t]{Q45 -- Explanation}
	\footnotesize
	\textbf{Correct: C.} Conclusions must be supported by sufficient evidence, and limits must be disclosed. \\[5pt]
\textbf{A wrong:} That issues an unsupported conclusion. \\[3pt]
\textbf{B wrong:} Boilerplate does not address lack of evidence. \\[3pt]
\textbf{D wrong:} Review is required. \\[3pt]
	\vfill {\tiny Set 2}
\end{frame}

% ===== Q46 =====
\begin{frame}[t]{Q46 \hfill {\scriptsize Ethics and Professionalism}}
	\small
	\textbf{Question:} While auditing inventory, an auditor finds indications of illegal waste dumping that are unrelated to her scope. What should she do?
	\begin{enumerate}[label=\Alph*.,itemsep=1pt,topsep=3pt]
			\item Ignore it, as it is out of scope.
			\item Report it promptly to the CAE so it can be communicated to appropriate parties, rather than ignoring it as out of scope.
			\item Confront the plant manager.
			\item Report only to a regulator.
	\end{enumerate}
\end{frame}

\begin{frame}[t]{Q46 -- Explanation}
	\footnotesize
	\textbf{Correct: B.} Auditors behave legally and ethically and communicate significant matters through the CAE. \\[5pt]
\textbf{A wrong:} Ignoring a possible legal violation is inappropriate. \\[3pt]
\textbf{C wrong:} Confrontation could compromise any inquiry. \\[3pt]
\textbf{D wrong:} Reporting outside the organization first bypasses governance channels. \\[3pt]
	\vfill {\tiny Set 2}
\end{frame}

% ===== Q47 =====
\begin{frame}[t]{Q47 \hfill {\scriptsize Ethics and Professionalism}}
	\small
	\textbf{Question:} An auditor is invited to join a management steering committee for a new system, with a vote. What is the BEST response?
	\begin{enumerate}[label=\Alph*.,itemsep=1pt,topsep=3pt]
			\item Decline voting membership; she may provide advice or observe while management retains decisions.
			\item Accept, since it builds knowledge.
			\item Accept but abstain from voting.
			\item Accept if membership is rotated annually.
	\end{enumerate}
\end{frame}

\begin{frame}[t]{Q47 -- Explanation}
	\footnotesize
	\textbf{Correct: A.} Participating in management decisions creates a self-review threat and can prevent later assurance. \\[5pt]
\textbf{B wrong:} Knowledge gain does not offset the impairment. \\[3pt]
\textbf{C wrong:} Membership still integrates her in management's decisions. \\[3pt]
\textbf{D wrong:} Rotation does not cure the problem. \\[3pt]
	\vfill {\tiny Set 2}
\end{frame}

% ===== Q48 =====
\begin{frame}[t]{Q48 \hfill {\scriptsize Ethics and Professionalism}}
	\small
	\textbf{Question:} Which safeguard BEST addresses a familiarity threat caused by long assignment to the same area?
	\begin{enumerate}[label=\Alph*.,itemsep=1pt,topsep=3pt]
			\item Increasing sample sizes.
			\item Using only a standard work program.
			\item Rotating assignments or adding independent review of the work.
			\item Shortening interview time.
	\end{enumerate}
\end{frame}

\begin{frame}[t]{Q48 -- Explanation}
	\footnotesize
	\textbf{Correct: C.} Rotation and independent review reduce the influence of long relationships. \\[5pt]
\textbf{A wrong:} Sample size does not address relationship bias. \\[3pt]
\textbf{B wrong:} A standard program does not address relationships. \\[3pt]
\textbf{D wrong:} Less interaction time does not fix the threat. \\[3pt]
	\vfill {\tiny Set 2}
\end{frame}

% ===== Q49 =====
\begin{frame}[t]{Q49 \hfill {\scriptsize Ethics and Professionalism}}
	\small
	\textbf{Question:} An auditor uses a machine-learning tool to flag anomalies but does not understand its limitations. Which principle is MOST at risk?
	\begin{enumerate}[label=\Alph*.,itemsep=1pt,topsep=3pt]
			\item Confidentiality.
			\item Competency.
			\item Organizational independence.
			\item Privilege against disclosure.
	\end{enumerate}
\end{frame}

\begin{frame}[t]{Q49 -- Explanation}
	\footnotesize
	\textbf{Correct: B.} Competency includes understanding tools and techniques used and their limitations. \\[5pt]
\textbf{A wrong:} Confidentiality is not the issue here. \\[3pt]
\textbf{C wrong:} Independence concerns reporting lines. \\[3pt]
\textbf{D wrong:} Privilege is a legal concept not relevant here. \\[3pt]
	\vfill {\tiny Set 2}
\end{frame}

% ===== Q50 =====
\begin{frame}[t]{Q50 \hfill {\scriptsize Ethics and Professionalism}}
	\small
	\textbf{Question:} An auditor learns of a pending acquisition during an engagement and mentions it to her spouse, who buys shares in the target. Which statement is MOST accurate?
	\begin{enumerate}[label=\Alph*.,itemsep=1pt,topsep=3pt]
			\item It is acceptable because the auditor did not trade herself.
			\item This violates confidentiality and integrity, because engagement information was used for personal gain.
			\item It is acceptable if the amount was small.
			\item It is acceptable if the deal is announced later.
	\end{enumerate}
\end{frame}

\begin{frame}[t]{Q50 -- Explanation}
	\footnotesize
	\textbf{Correct: B.} Engagement information must not be used for personal gain or disclosed improperly, directly or through others. \\[5pt]
\textbf{A wrong:} Enabling a family member to profit still misuses the information. \\[3pt]
\textbf{C wrong:} Size does not remove the violation. \\[3pt]
\textbf{D wrong:} Later announcement does not cure the misuse. \\[3pt]
	\vfill {\tiny Set 2}
\end{frame}

% ===== Q51 =====
\begin{frame}[t]{Q51 \hfill {\scriptsize Ethics and Professionalism}}
	\small
	\textbf{Question:} In exercising due professional care, which factor is LEAST relevant?
	\begin{enumerate}[label=\Alph*.,itemsep=1pt,topsep=3pt]
			\item The probability of significant errors, fraud, or noncompliance.
			\item The seniority of the auditee manager.
			\item The complexity and extent of work needed.
			\item The cost of assurance relative to its potential benefits.
	\end{enumerate}
\end{frame}

\begin{frame}[t]{Q51 -- Explanation}
	\footnotesize
	\textbf{Correct: B.} Due care weighs risk, complexity, and cost-benefit; the auditee's rank is not a factor in procedures needed. \\[5pt]
\textbf{A wrong:} Relevant to the work needed. \\[3pt]
\textbf{C wrong:} Relevant to the work needed. \\[3pt]
\textbf{D wrong:} Relevant to the work needed. \\[3pt]
	\vfill {\tiny Set 2}
\end{frame}

% ===== Q52 =====
\begin{frame}[t]{Q52 \hfill {\scriptsize Ethics and Professionalism}}
	\small
	\textbf{Question:} An auditee manager asks the CAE to remove the auditor leading the engagement after she raised findings. What is the MOST appropriate response?
	\begin{enumerate}[label=\Alph*.,itemsep=1pt,topsep=3pt]
			\item Evaluate any substantiated concerns about conduct objectively, but decline reassignment solely because the auditee dislikes the findings.
			\item Reassign her to keep relations good.
			\item Discard the findings.
			\item Let the manager choose the replacement.
	\end{enumerate}
\end{frame}

\begin{frame}[t]{Q52 -- Explanation}
	\footnotesize
	\textbf{Correct: A.} The auditee should not control audit assignments or outcomes; legitimate conduct concerns are still reviewed. \\[5pt]
\textbf{B wrong:} That lets the auditee influence the work. \\[3pt]
\textbf{C wrong:} Findings must stand on evidence. \\[3pt]
\textbf{D wrong:} The auditee's choice would impair independence. \\[3pt]
	\vfill {\tiny Set 2}
\end{frame}

% ===== Q53 =====
\begin{frame}[t]{Q53 \hfill {\scriptsize Ethics and Professionalism}}
	\small
	\textbf{Question:} An auditee offers an auditor a \$300 sports ticket a few days before the final report. What should the auditor do?
	\begin{enumerate}[label=\Alph*.,itemsep=1pt,topsep=3pt]
			\item Accept it if it is within the auditee's gift policy.
			\item Accept it and disclose it after the report.
			\item Accept it and donate it to charity.
			\item Decline it and inform the CAE, because it could reasonably be seen as impairing objectivity.
	\end{enumerate}
\end{frame}

\begin{frame}[t]{Q53 -- Explanation}
	\footnotesize
	\textbf{Correct: D.} A gift of this value from an auditee during an engagement creates an appearance of impaired objectivity. \\[5pt]
\textbf{A wrong:} The auditee's policy does not govern the auditor's obligations. \\[3pt]
\textbf{B wrong:} Late disclosure does not remove the threat. \\[3pt]
\textbf{C wrong:} Donating does not change that it was accepted. \\[3pt]
	\vfill {\tiny Set 2}
\end{frame}

% ===== Q54 =====
\begin{frame}[t]{Q54 \hfill {\scriptsize Ethics and Professionalism}}
	\small
	\textbf{Question:} An auditor uses the organization's audit methodology and templates in a private consulting engagement. Which statement is MOST accurate?
	\begin{enumerate}[label=\Alph*.,itemsep=1pt,topsep=3pt]
			\item It is acceptable since the auditor developed some of the templates.
			\item It is acceptable if clients are not competitors.
			\item This misuses information and materials acquired through employment, contrary to confidentiality and integrity.
			\item It is acceptable if unpaid.
	\end{enumerate}
\end{frame}

\begin{frame}[t]{Q54 -- Explanation}
	\footnotesize
	\textbf{Correct: C.} Organizational materials are not for personal use without authority. \\[5pt]
\textbf{A wrong:} Authorship in the course of employment does not transfer ownership. \\[3pt]
\textbf{B wrong:} Client identity does not matter. \\[3pt]
\textbf{D wrong:} Payment is not the issue. \\[3pt]
	\vfill {\tiny Set 2}
\end{frame}

% ===== Q55 =====
\begin{frame}[t]{Q55 \hfill {\scriptsize Ethics and Professionalism}}
	\small
	\textbf{Question:} The audit committee chair privately asks the CAE to downplay a finding concerning a company owned by the chair's relative. What should the CAE do?
	\begin{enumerate}[label=\Alph*.,itemsep=1pt,topsep=3pt]
			\item Downplay the finding to keep the chair's support.
			\item Quietly remove the finding.
			\item Report the finding accurately and raise the request with the full board or governing body.
			\item Tell only the external auditor.
	\end{enumerate}
\end{frame}

\begin{frame}[t]{Q55 -- Explanation}
	\footnotesize
	\textbf{Correct: C.} Integrity and professional courage require accurate reporting despite pressure, and the request should be escalated. \\[5pt]
\textbf{A wrong:} Downplaying distorts the message. \\[3pt]
\textbf{B wrong:} Removing a valid finding misleads users. \\[3pt]
\textbf{D wrong:} The external auditor does not resolve the internal reporting problem. \\[3pt]
	\vfill {\tiny Set 2}
\end{frame}

% ===== Q56 =====
\begin{frame}[t]{Q56 \hfill {\scriptsize Governance, Risk Management and Control}}
	\small
	\textbf{Question:} Which is a DIRECTIVE control?
	\begin{enumerate}[label=\Alph*.,itemsep=1pt,topsep=3pt]
			\item A written policy requiring dual authorization for payments above a set limit.
			\item A locked server room.
			\item A weekly exception report review.
			\item Restoring from backups after an outage.
	\end{enumerate}
\end{frame}

\begin{frame}[t]{Q56 -- Explanation}
	\footnotesize
	\textbf{Correct: A.} Directive controls such as policies and procedures guide behavior toward desired outcomes. \\[5pt]
\textbf{B wrong:} That is preventive. \\[3pt]
\textbf{C wrong:} That is detective. \\[3pt]
\textbf{D wrong:} That is corrective. \\[3pt]
	\vfill {\tiny Set 2}
\end{frame}

% ===== Q57 =====
\begin{frame}[t]{Q57 \hfill {\scriptsize Governance, Risk Management and Control}}
	\small
	\textbf{Question:} Which factor MOST affects the effectiveness of a management review control?
	\begin{enumerate}[label=\Alph*.,itemsep=1pt,topsep=3pt]
			\item The seniority of the reviewer alone.
			\item The precision of the review, including thresholds and follow-up of exceptions.
			\item The length of the review report.
			\item Whether the review is documented only by email.
	\end{enumerate}
\end{frame}

\begin{frame}[t]{Q57 -- Explanation}
	\footnotesize
	\textbf{Correct: B.} A review control works if it is precise enough to catch errors and exceptions are followed up. \\[5pt]
\textbf{A wrong:} Seniority does not ensure precision. \\[3pt]
\textbf{C wrong:} Length does not ensure quality. \\[3pt]
\textbf{D wrong:} The medium of documentation does not affect operation. \\[3pt]
	\vfill {\tiny Set 2}
\end{frame}

% ===== Q58 =====
\begin{frame}[t]{Q58 \hfill {\scriptsize Governance, Risk Management and Control}}
	\small
	\textbf{Question:} Under COSO's 2013 Internal Control - Integrated Framework, the principle of commitment to integrity and ethical values belongs to which component?
	\begin{enumerate}[label=\Alph*.,itemsep=1pt,topsep=3pt]
			\item Control Environment.
			\item Control Activities.
			\item Information and Communication.
			\item Monitoring Activities.
	\end{enumerate}
\end{frame}

\begin{frame}[t]{Q58 -- Explanation}
	\footnotesize
	\textbf{Correct: A.} The control environment includes commitment to integrity and ethical values, oversight, structure, competence, and accountability. \\[5pt]
\textbf{B wrong:} Control activities are the actions that mitigate risks. \\[3pt]
\textbf{C wrong:} That component covers information quality and communication. \\[3pt]
\textbf{D wrong:} Monitoring evaluates the system of internal control. \\[3pt]
	\vfill {\tiny Set 2}
\end{frame}

% ===== Q59 =====
\begin{frame}[t]{Q59 \hfill {\scriptsize Governance, Risk Management and Control}}
	\small
	\textbf{Question:} Ransomware is detected spreading on a file server. What should be the FIRST priority?
	\begin{enumerate}[label=\Alph*.,itemsep=1pt,topsep=3pt]
			\item Pay the ransom immediately.
			\item Restore from backups before isolating systems.
			\item Wipe affected servers to remove the malware.
			\item Contain the incident by isolating affected systems while preserving evidence.
	\end{enumerate}
\end{frame}

\begin{frame}[t]{Q59 -- Explanation}
	\footnotesize
	\textbf{Correct: D.} Containment limits spread; evidence preservation supports investigation. \\[5pt]
\textbf{A wrong:} Payment is not a first step and does not guarantee recovery. \\[3pt]
\textbf{B wrong:} Restoring before isolation risks reinfection. \\[3pt]
\textbf{C wrong:} Wiping destroys evidence and may not stop spread. \\[3pt]
	\vfill {\tiny Set 2}
\end{frame}

% ===== Q60 =====
\begin{frame}[t]{Q60 \hfill {\scriptsize Governance, Risk Management and Control}}
	\small
	\textbf{Question:} Which situation MOST indicates that risk appetite is not embedded in decision-making?
	\begin{enumerate}[label=\Alph*.,itemsep=1pt,topsep=3pt]
			\item Limits are reviewed annually by the board.
			\item Exceptions are escalated and logged.
			\item Business units approve exposures beyond approved limits without escalation.
			\item Appetite statements are linked to strategic objectives.
	\end{enumerate}
\end{frame}

\begin{frame}[t]{Q60 -- Explanation}
	\footnotesize
	\textbf{Correct: C.} Embedding means decisions are made within limits, with escalation of exceptions. \\[5pt]
\textbf{A wrong:} That is a sound practice. \\[3pt]
\textbf{B wrong:} That shows escalation works. \\[3pt]
\textbf{D wrong:} That supports embedding. \\[3pt]
	\vfill {\tiny Set 2}
\end{frame}

% ===== Q61 =====
\begin{frame}[t]{Q61 \hfill {\scriptsize Governance, Risk Management and Control}}
	\small
	\textbf{Question:} Which statement BEST describes a key risk indicator (KRI)?
	\begin{enumerate}[label=\Alph*.,itemsep=1pt,topsep=3pt]
			\item A metric that gives early warning of increasing risk exposure, such as rising turnover in a critical process.
			\item A measure of losses already incurred only.
			\item The annual external audit opinion.
			\item The list of approved controls.
	\end{enumerate}
\end{frame}

\begin{frame}[t]{Q61 -- Explanation}
	\footnotesize
	\textbf{Correct: A.} KRIs are generally forward-looking metrics that signal changes in risk exposure. \\[5pt]
\textbf{B wrong:} That describes a lagging metric. \\[3pt]
\textbf{C wrong:} An audit opinion is not an indicator. \\[3pt]
\textbf{D wrong:} A control list is not a metric. \\[3pt]
	\vfill {\tiny Set 2}
\end{frame}

% ===== Q62 =====
\begin{frame}[t]{Q62 \hfill {\scriptsize Governance, Risk Management and Control}}
	\small
	\textbf{Question:} Which control BEST ensures complete and accurate transfer of transactions between two systems?
	\begin{enumerate}[label=\Alph*.,itemsep=1pt,topsep=3pt]
			\item A quarterly interview with the interface owner.
			\item Automated reconciliation of record counts and control totals with exception follow-up.
			\item A screenshot of the interface status.
			\item An annual review of the interface specification.
	\end{enumerate}
\end{frame}

\begin{frame}[t]{Q62 -- Explanation}
	\footnotesize
	\textbf{Correct: B.} Automated count and total reconciliation detects missing or altered records promptly. \\[5pt]
\textbf{A wrong:} Interviews do not test data transfer. \\[3pt]
\textbf{C wrong:} A screenshot shows status only. \\[3pt]
\textbf{D wrong:} Specification review is infrequent and design-only. \\[3pt]
	\vfill {\tiny Set 2}
\end{frame}

% ===== Q63 =====
\begin{frame}[t]{Q63 \hfill {\scriptsize Governance, Risk Management and Control}}
	\small
	\textbf{Question:} A single loss event is estimated at \$200,000 with an annual rate of occurrence of 0.25. A control reduces the rate to 0.05 and costs \$30,000 per year. What is the net annual benefit?
	\begin{enumerate}[label=\Alph*.,itemsep=1pt,topsep=3pt]
			\item \$10,000.
			\item \$40,000.
			\item \$20,000.
			\item \$70,000.
	\end{enumerate}
\end{frame}

\begin{frame}[t]{Q63 -- Explanation}
	\footnotesize
	\textbf{Correct: A.} Annual loss expectancy falls from \$50,000 to \$10,000, a \$40,000 reduction. Less the \$30,000 cost, the net benefit is \$10,000. \\[5pt]
\textbf{B wrong:} That is the gross reduction before cost. \\[3pt]
\textbf{C wrong:} That subtracts the cost from the original expectancy rather than from the reduction. \\[3pt]
\textbf{D wrong:} That adds the cost instead of subtracting it. \\[3pt]
	\vfill {\tiny Set 2}
\end{frame}

% ===== Q64 =====
\begin{frame}[t]{Q64 \hfill {\scriptsize Governance, Risk Management and Control}}
	\small
	\textbf{Question:} In the COSO ERM framework, which component addresses identifying, assessing, and prioritizing risks and selecting responses?
	\begin{enumerate}[label=\Alph*.,itemsep=1pt,topsep=3pt]
			\item Governance and Culture.
			\item Strategy and Objective-Setting.
			\item Review and Revision.
			\item Performance.
	\end{enumerate}
\end{frame}

\begin{frame}[t]{Q64 -- Explanation}
	\footnotesize
	\textbf{Correct: D.} The Performance component covers risk identification, assessment, prioritization, response, and portfolio view. \\[5pt]
\textbf{A wrong:} It covers oversight, culture, and values. \\[3pt]
\textbf{B wrong:} It covers strategy and objectives. \\[3pt]
\textbf{C wrong:} It covers revising practices over time. \\[3pt]
	\vfill {\tiny Set 2}
\end{frame}

% ===== Q65 =====
\begin{frame}[t]{Q65 \hfill {\scriptsize Governance, Risk Management and Control}}
	\small
	\textbf{Question:} Which provides the strongest evidence that a vulnerability is actually exploitable?
	\begin{enumerate}[label=\Alph*.,itemsep=1pt,topsep=3pt]
			\item An automated vulnerability scan.
			\item A policy review.
			\item An asset inventory.
			\item A penetration test.
	\end{enumerate}
\end{frame}

\begin{frame}[t]{Q65 -- Explanation}
	\footnotesize
	\textbf{Correct: D.} Penetration testing attempts exploitation and demonstrates real impact. \\[5pt]
\textbf{A wrong:} Scans identify potential weaknesses, not proven exploitability. \\[3pt]
\textbf{B wrong:} Policies do not test exploitability. \\[3pt]
\textbf{C wrong:} An inventory lists assets only. \\[3pt]
	\vfill {\tiny Set 2}
\end{frame}

% ===== Q66 =====
\begin{frame}[t]{Q66 \hfill {\scriptsize Governance, Risk Management and Control}}
	\small
	\textbf{Question:} Which characteristic MOST determines whether a detective control is effective?
	\begin{enumerate}[label=\Alph*.,itemsep=1pt,topsep=3pt]
			\item The number of exceptions generated.
			\item Whether it is documented in a policy.
			\item Timely investigation and resolution of the exceptions it identifies.
			\item Whether the reviewer is senior.
	\end{enumerate}
\end{frame}

\begin{frame}[t]{Q66 -- Explanation}
	\footnotesize
	\textbf{Correct: C.} A detective control works only if exceptions are followed up and corrected. \\[5pt]
\textbf{A wrong:} Volume does not show effectiveness. \\[3pt]
\textbf{B wrong:} Documentation alone does not prove operation. \\[3pt]
\textbf{D wrong:} Seniority does not ensure follow-up. \\[3pt]
	\vfill {\tiny Set 2}
\end{frame}

% ===== Q67 =====
\begin{frame}[t]{Q67 \hfill {\scriptsize Governance, Risk Management and Control}}
	\small
	\textbf{Question:} Which control BEST maintains segregation of duties in a CI/CD environment?
	\begin{enumerate}[label=\Alph*.,itemsep=1pt,topsep=3pt]
			\item Developers approving their own pull requests.
			\item Shared deployment credentials known to the team.
			\item Post-deployment review once a quarter.
			\item Pipeline-enforced peer review and independent approval before deployment, with developers unable to push directly to production.
	\end{enumerate}
\end{frame}

\begin{frame}[t]{Q67 -- Explanation}
	\footnotesize
	\textbf{Correct: D.} Enforced review and approval gates prevent unauthorized or untested changes. \\[5pt]
\textbf{A wrong:} Self-approval removes independence. \\[3pt]
\textbf{B wrong:} Shared credentials remove accountability. \\[3pt]
\textbf{C wrong:} Infrequent review is too late. \\[3pt]
	\vfill {\tiny Set 2}
\end{frame}

% ===== Q68 =====
\begin{frame}[t]{Q68 \hfill {\scriptsize Governance, Risk Management and Control}}
	\small
	\textbf{Question:} Which control is MOST important for emergency (firefighter) privileged access?
	\begin{enumerate}[label=\Alph*.,itemsep=1pt,topsep=3pt]
			\item A shared password kept in a team spreadsheet.
			\item Permanent privileged access for the on-call team.
			\item Time-limited access with approval, full logging, and independent review after use.
			\item Verbal approval with no record.
	\end{enumerate}
\end{frame}

\begin{frame}[t]{Q68 -- Explanation}
	\footnotesize
	\textbf{Correct: C.} Emergency access needs approval, limits, logging, and after-the-fact review. \\[5pt]
\textbf{A wrong:} Shared passwords destroy accountability. \\[3pt]
\textbf{B wrong:} Permanent privileged access defeats least privilege. \\[3pt]
\textbf{D wrong:} No record prevents review. \\[3pt]
	\vfill {\tiny Set 2}
\end{frame}

% ===== Q69 =====
\begin{frame}[t]{Q69 \hfill {\scriptsize Governance, Risk Management and Control}}
	\small
	\textbf{Question:} A critical cloud vendor provides no independent assurance report. What should internal audit recommend?
	\begin{enumerate}[label=\Alph*.,itemsep=1pt,topsep=3pt]
			\item Use contractual audit rights or other assurance, such as questionnaires and independent testing, and assess the risk of reliance.
			\item Accept the vendor's marketing statements.
			\item Assume risk is transferred to the vendor.
			\item Stop all use immediately regardless of impact.
	\end{enumerate}
\end{frame}

\begin{frame}[t]{Q69 -- Explanation}
	\footnotesize
	\textbf{Correct: A.} Outsourcing does not transfer accountability; the organization needs alternative assurance. \\[5pt]
\textbf{B wrong:} Marketing is not evidence. \\[3pt]
\textbf{C wrong:} Accountability remains with the organization. \\[3pt]
\textbf{D wrong:} Immediate termination ignores business impact and may not be necessary. \\[3pt]
	\vfill {\tiny Set 2}
\end{frame}

% ===== Q70 =====
\begin{frame}[t]{Q70 \hfill {\scriptsize Governance, Risk Management and Control}}
	\small
	\textbf{Question:} Which practice BEST supports accountability for AI use across the organization?
	\begin{enumerate}[label=\Alph*.,itemsep=1pt,topsep=3pt]
			\item A policy statement with no register.
			\item Allowing each team to adopt tools independently.
			\item An inventory of AI systems with named owners, risk classification, and approvals at key lifecycle stages.
			\item Reviewing vendor brochures annually.
	\end{enumerate}
\end{frame}

\begin{frame}[t]{Q70 -- Explanation}
	\footnotesize
	\textbf{Correct: C.} An inventory with ownership and lifecycle approvals makes AI use visible and governed. \\[5pt]
\textbf{A wrong:} Without a register, use is not visible. \\[3pt]
\textbf{B wrong:} Uncoordinated adoption weakens governance. \\[3pt]
\textbf{D wrong:} Brochures are not governance. \\[3pt]
	\vfill {\tiny Set 2}
\end{frame}

% ===== Q71 =====
\begin{frame}[t]{Q71 \hfill {\scriptsize Governance, Risk Management and Control}}
	\small
	\textbf{Question:} Under GDPR, how quickly must a controller notify the supervisory authority of a personal data breach, where feasible?
	\begin{enumerate}[label=\Alph*.,itemsep=1pt,topsep=3pt]
			\item Within 24 hours of the breach occurring.
			\item Within 30 days of discovery.
			\item Without undue delay and, where feasible, within 72 hours of becoming aware.
			\item Only if the breach affects more than 10,000 individuals.
	\end{enumerate}
\end{frame}

\begin{frame}[t]{Q71 -- Explanation}
	\footnotesize
	\textbf{Correct: C.} The GDPR sets a 72-hour expectation from awareness, unless the breach is unlikely to result in risk. \\[5pt]
\textbf{A wrong:} 24 hours is not the requirement. \\[3pt]
\textbf{B wrong:} 30 days is too long. \\[3pt]
\textbf{D wrong:} There is no fixed individual-count threshold. \\[3pt]
	\vfill {\tiny Set 2}
\end{frame}

% ===== Q72 =====
\begin{frame}[t]{Q72 \hfill {\scriptsize Governance, Risk Management and Control}}
	\small
	\textbf{Question:} Which statement BEST reflects a zero trust approach?
	\begin{enumerate}[label=\Alph*.,itemsep=1pt,topsep=3pt]
			\item Anything inside the corporate network is trusted.
			\item Strong perimeter firewalls remove the need for user authentication.
			\item Once authenticated, users keep access until their account is closed.
			\item No user or device is trusted by default; access is verified continuously and limited by least privilege.
	\end{enumerate}
\end{frame}

\begin{frame}[t]{Q72 -- Explanation}
	\footnotesize
	\textbf{Correct: D.} Zero trust assumes no implicit trust and verifies each request. \\[5pt]
\textbf{A wrong:} That is the perimeter model. \\[3pt]
\textbf{B wrong:} Firewalls do not replace authentication. \\[3pt]
\textbf{C wrong:} Persistent access contradicts continuous verification. \\[3pt]
	\vfill {\tiny Set 2}
\end{frame}

% ===== Q73 =====
\begin{frame}[t]{Q73 \hfill {\scriptsize Governance, Risk Management and Control}}
	\small
	\textbf{Question:} Which control BEST addresses the risk of AI model drift?
	\begin{enumerate}[label=\Alph*.,itemsep=1pt,topsep=3pt]
			\item Ongoing performance monitoring against baselines with thresholds that trigger review or retraining.
			\item Annual user training.
			\item One-time validation before go-live.
			\item Restricting access to the model's source code.
	\end{enumerate}
\end{frame}

\begin{frame}[t]{Q73 -- Explanation}
	\footnotesize
	\textbf{Correct: A.} Drift develops over time, so continuous monitoring and defined triggers are needed. \\[5pt]
\textbf{B wrong:} Training does not detect drift. \\[3pt]
\textbf{C wrong:} A one-time check cannot detect later changes. \\[3pt]
\textbf{D wrong:} Code access does not address performance changes. \\[3pt]
	\vfill {\tiny Set 2}
\end{frame}

% ===== Q74 =====
\begin{frame}[t]{Q74 \hfill {\scriptsize Governance, Risk Management and Control}}
	\small
	\textbf{Question:} Which provides the greatest assurance that recovery of a critical system will work?
	\begin{enumerate}[label=\Alph*.,itemsep=1pt,topsep=3pt]
			\item Reviewing the continuity plan document.
			\item A tabletop discussion of recovery steps.
			\item A live failover test of the system at the recovery site.
			\item Confirming that staff have read the plan.
	\end{enumerate}
\end{frame}

\begin{frame}[t]{Q74 -- Explanation}
	\footnotesize
	\textbf{Correct: C.} Actual failover testing demonstrates that recovery capabilities work in practice. \\[5pt]
\textbf{A wrong:} Document review does not test capability. \\[3pt]
\textbf{B wrong:} Tabletop exercises test understanding, not technology. \\[3pt]
\textbf{D wrong:} Awareness does not prove recoverability. \\[3pt]
	\vfill {\tiny Set 2}
\end{frame}

% ===== Q75 =====
\begin{frame}[t]{Q75 \hfill {\scriptsize Governance, Risk Management and Control}}
	\small
	\textbf{Question:} For a risk with low likelihood but severe impact, which response is typically MOST appropriate?
	\begin{enumerate}[label=\Alph*.,itemsep=1pt,topsep=3pt]
			\item Accept it with no further action.
			\item Avoid it by exiting all operations.
			\item Rely on historical frequency to ignore it.
			\item Share part of the risk, for example through insurance, and prepare contingency plans.
	\end{enumerate}
\end{frame}

\begin{frame}[t]{Q75 -- Explanation}
	\footnotesize
	\textbf{Correct: D.} Low-frequency, high-severity risks are often addressed by sharing and contingency planning. \\[5pt]
\textbf{A wrong:} Unmonitored acceptance ignores severe impact. \\[3pt]
\textbf{B wrong:} Exiting operations is disproportionate. \\[3pt]
\textbf{C wrong:} Historical frequency can understate rare events. \\[3pt]
	\vfill {\tiny Set 2}
\end{frame}

% ===== Q76 =====
\begin{frame}[t]{Q76 \hfill {\scriptsize Governance, Risk Management and Control}}
	\small
	\textbf{Question:} Which is the MOST significant finding in an audit of robotic process automation (RPA)?
	\begin{enumerate}[label=\Alph*.,itemsep=1pt,topsep=3pt]
			\item Bots run under shared privileged credentials with no named owner or logging.
			\item Bot documentation uses an old template.
			\item Bot schedules are reviewed monthly by owners.
			\item Bot IDs are unique and have least-privilege access.
	\end{enumerate}
\end{frame}

\begin{frame}[t]{Q76 -- Explanation}
	\footnotesize
	\textbf{Correct: A.} Shared privileged credentials without ownership or logging create accountability and security gaps. \\[5pt]
\textbf{B wrong:} A template issue is minor. \\[3pt]
\textbf{C wrong:} Owner review is a good practice. \\[3pt]
\textbf{D wrong:} Unique least-privilege IDs are a good control. \\[3pt]
	\vfill {\tiny Set 2}
\end{frame}

% ===== Q77 =====
\begin{frame}[t]{Q77 \hfill {\scriptsize Governance, Risk Management and Control}}
	\small
	\textbf{Question:} Which is an ENTITY-LEVEL control?
	\begin{enumerate}[label=\Alph*.,itemsep=1pt,topsep=3pt]
			\item A three-way match on supplier invoices.
			\item An edit check on invoice dates.
			\item A daily bank reconciliation for one account.
			\item A whistleblower hotline and enforcement of the code of conduct.
	\end{enumerate}
\end{frame}

\begin{frame}[t]{Q77 -- Explanation}
	\footnotesize
	\textbf{Correct: D.} Entity-level controls such as ethics programs operate across the organization. \\[5pt]
\textbf{A wrong:} That is a transaction-level control. \\[3pt]
\textbf{B wrong:} That is an application control. \\[3pt]
\textbf{C wrong:} That is an account-level activity. \\[3pt]
	\vfill {\tiny Set 2}
\end{frame}

% ===== Q78 =====
\begin{frame}[t]{Q78 \hfill {\scriptsize Governance, Risk Management and Control}}
	\small
	\textbf{Question:} Which is the BEST indicator of a healthy risk culture?
	\begin{enumerate}[label=\Alph*.,itemsep=1pt,topsep=3pt]
			\item Zero reported incidents over several years.
			\item A detailed risk policy exists.
			\item Employees feel safe raising concerns and leaders act on them promptly.
			\item Risk reports are issued monthly.
	\end{enumerate}
\end{frame}

\begin{frame}[t]{Q78 -- Explanation}
	\footnotesize
	\textbf{Correct: C.} Healthy culture shows in behavior: open escalation and leadership response. \\[5pt]
\textbf{A wrong:} Zero incidents may indicate under-reporting. \\[3pt]
\textbf{B wrong:} A policy alone does not show behavior. \\[3pt]
\textbf{D wrong:} Reporting frequency does not show culture. \\[3pt]
	\vfill {\tiny Set 2}
\end{frame}

% ===== Q79 =====
\begin{frame}[t]{Q79 \hfill {\scriptsize Governance, Risk Management and Control}}
	\small
	\textbf{Question:} Inherent risk is rated impact 5 and likelihood 4 (score 20). Controls reduce likelihood to 2 and impact is unchanged. What is the residual score?
	\begin{enumerate}[label=\Alph*.,itemsep=1pt,topsep=3pt]
			\item 8.
			\item 15.
			\item 10.
			\item 20.
	\end{enumerate}
\end{frame}

\begin{frame}[t]{Q79 -- Explanation}
	\footnotesize
	\textbf{Correct: C.} Residual score is 5 x 2 = 10. \\[5pt]
\textbf{A wrong:} That would need impact of 4. \\[3pt]
\textbf{B wrong:} That would reduce likelihood only to 3. \\[3pt]
\textbf{D wrong:} That ignores the effect of controls. \\[3pt]
	\vfill {\tiny Set 2}
\end{frame}

% ===== Q80 =====
\begin{frame}[t]{Q80 \hfill {\scriptsize Governance, Risk Management and Control}}
	\small
	\textbf{Question:} Which ERM component addresses evaluating alternative strategies and the risk profile they imply?
	\begin{enumerate}[label=\Alph*.,itemsep=1pt,topsep=3pt]
			\item Review and Revision.
			\item Strategy and Objective-Setting.
			\item Information, Communication, and Reporting.
			\item Governance and Culture.
	\end{enumerate}
\end{frame}

\begin{frame}[t]{Q80 -- Explanation}
	\footnotesize
	\textbf{Correct: B.} Strategy and Objective-Setting links risk appetite, alternatives, and business objectives. \\[5pt]
\textbf{A wrong:} It addresses changes and improvements over time. \\[3pt]
\textbf{C wrong:} It addresses information flow. \\[3pt]
\textbf{D wrong:} It addresses oversight and culture. \\[3pt]
	\vfill {\tiny Set 2}
\end{frame}

% ===== Q81 =====
\begin{frame}[t]{Q81 \hfill {\scriptsize Governance, Risk Management and Control}}
	\small
	\textbf{Question:} Which cloud finding is MOST significant?
	\begin{enumerate}[label=\Alph*.,itemsep=1pt,topsep=3pt]
			\item Resource tags are inconsistent.
			\item The root (account owner) credentials are used daily without multi-factor authentication.
			\item Some test instances are left running.
			\item Billing alerts are set above the policy threshold.
	\end{enumerate}
\end{frame}

\begin{frame}[t]{Q81 -- Explanation}
	\footnotesize
	\textbf{Correct: B.} Unprotected, routinely used root credentials expose the entire cloud environment. \\[5pt]
\textbf{A wrong:} Tagging is an administrative issue. \\[3pt]
\textbf{C wrong:} Idle instances are mainly a cost issue. \\[3pt]
\textbf{D wrong:} Alert thresholds are a minor monitoring issue. \\[3pt]
	\vfill {\tiny Set 2}
\end{frame}

% ===== Q82 =====
\begin{frame}[t]{Q82 \hfill {\scriptsize Governance, Risk Management and Control}}
	\small
	\textbf{Question:} Which control BEST protects the integrity of system logs?
	\begin{enumerate}[label=\Alph*.,itemsep=1pt,topsep=3pt]
			\item Storing logs on the same server with administrator write access.
			\item Allowing application owners to edit logs for clarity.
			\item Deleting logs after seven days to save space.
			\item Sending logs to a centralized, write-once repository with restricted access and monitoring.
	\end{enumerate}
\end{frame}

\begin{frame}[t]{Q82 -- Explanation}
	\footnotesize
	\textbf{Correct: D.} Centralized, tamper-resistant storage preserves logs for investigation. \\[5pt]
\textbf{A wrong:} Local storage with admin write access can be altered. \\[3pt]
\textbf{B wrong:} Editable logs lose evidentiary value. \\[3pt]
\textbf{C wrong:} Short retention loses evidence. \\[3pt]
	\vfill {\tiny Set 2}
\end{frame}

% ===== Q83 =====
\begin{frame}[t]{Q83 \hfill {\scriptsize Governance, Risk Management and Control}}
	\small
	\textbf{Question:} Employees say middle managers routinely bypass the code of conduct despite senior leaders' statements. What is the MOST appropriate internal audit response?
	\begin{enumerate}[label=\Alph*.,itemsep=1pt,topsep=3pt]
			\item Disregard it because senior leaders set the tone.
			\item Treat it as an HR matter outside internal audit's scope.
			\item Assess the cultural risk, report it to senior management and the board, and recommend actions to align behavior with stated values.
			\item Issue findings against named managers without further work.
	\end{enumerate}
\end{frame}

\begin{frame}[t]{Q83 -- Explanation}
	\footnotesize
	\textbf{Correct: C.} Culture is part of governance and control environment, and internal audit can assess and report on it. \\[5pt]
\textbf{A wrong:} Tone in the middle also shapes behavior. \\[3pt]
\textbf{B wrong:} Ethics culture is within internal audit's scope. \\[3pt]
\textbf{D wrong:} Findings need evidence and due process. \\[3pt]
	\vfill {\tiny Set 2}
\end{frame}

% ===== Q84 =====
\begin{frame}[t]{Q84 \hfill {\scriptsize Governance, Risk Management and Control}}
	\small
	\textbf{Question:} Who should OWN a specific operational risk?
	\begin{enumerate}[label=\Alph*.,itemsep=1pt,topsep=3pt]
			\item The manager accountable for the objective the risk affects.
			\item The head of the risk function.
			\item The CAE.
			\item The external auditor.
	\end{enumerate}
\end{frame}

\begin{frame}[t]{Q84 -- Explanation}
	\footnotesize
	\textbf{Correct: A.} Risk ownership belongs with those accountable for the related objectives (first line). \\[5pt]
\textbf{B wrong:} The risk function advises and monitors. \\[3pt]
\textbf{C wrong:} Internal audit provides independent assurance. \\[3pt]
\textbf{D wrong:} External auditors have no ownership role. \\[3pt]
	\vfill {\tiny Set 2}
\end{frame}

% ===== Q85 =====
\begin{frame}[t]{Q85 \hfill {\scriptsize Governance, Risk Management and Control}}
	\small
	\textbf{Question:} An enterprise chatbot answers questions using internal documents. Which control BEST reduces the risk of exposing sensitive information to unauthorized users?
	\begin{enumerate}[label=\Alph*.,itemsep=1pt,topsep=3pt]
			\item Telling users not to ask sensitive questions.
			\item Enforcing source-system access permissions on retrieval, with data loss prevention and logging.
			\item Indexing all documents so answers are complete.
			\item Relying on a longer system prompt only.
	\end{enumerate}
\end{frame}

\begin{frame}[t]{Q85 -- Explanation}
	\footnotesize
	\textbf{Correct: B.} Technical access enforcement and monitoring are needed rather than reliance on user behavior or prompts. \\[5pt]
\textbf{A wrong:} User instructions are not a control. \\[3pt]
\textbf{C wrong:} Indexing everything increases exposure. \\[3pt]
\textbf{D wrong:} Prompts alone can be bypassed. \\[3pt]
	\vfill {\tiny Set 2}
\end{frame}

% ===== Q86 =====
\begin{frame}[t]{Q86 \hfill {\scriptsize Fraud Risks}}
	\small
	\textbf{Question:} Which technique is MOST commonly used to conceal financial statement fraud by senior management?
	\begin{enumerate}[label=\Alph*.,itemsep=1pt,topsep=3pt]
			\item Period-end top-side journal entries without supporting documentation.
			\item Small expense claims.
			\item Employee cash theft.
			\item Vendor discounts.
	\end{enumerate}
\end{frame}

\begin{frame}[t]{Q86 -- Explanation}
	\footnotesize
	\textbf{Correct: A.} Senior management can post unsupported period-end entries that bypass normal controls. \\[5pt]
\textbf{B wrong:} That is a small-scale asset misappropriation. \\[3pt]
\textbf{C wrong:} That is an employee asset scheme. \\[3pt]
\textbf{D wrong:} Discounts are routine. \\[3pt]
	\vfill {\tiny Set 2}
\end{frame}

% ===== Q87 =====
\begin{frame}[t]{Q87 \hfill {\scriptsize Fraud Risks}}
	\small
	\textbf{Question:} Accounts payable receives an urgent email appearing to be from the CEO instructing a change to a vendor's bank details. Which control BEST addresses this risk?
	\begin{enumerate}[label=\Alph*.,itemsep=1pt,topsep=3pt]
			\item Replying to the email to confirm.
			\item Call-back verification using an independently sourced phone number, and dual approval of bank detail changes.
			\item Increasing email storage capacity.
			\item Requiring an annual password change.
	\end{enumerate}
\end{frame}

\begin{frame}[t]{Q87 -- Explanation}
	\footnotesize
	\textbf{Correct: B.} Out-of-band verification and dual approval counter business email compromise. \\[5pt]
\textbf{A wrong:} Replying uses the possibly compromised channel. \\[3pt]
\textbf{C wrong:} Storage is unrelated. \\[3pt]
\textbf{D wrong:} Password changes do not stop impersonation. \\[3pt]
	\vfill {\tiny Set 2}
\end{frame}

% ===== Q88 =====
\begin{frame}[t]{Q88 \hfill {\scriptsize Fraud Risks}}
	\small
	\textbf{Question:} Under pressure to meet targets, a sales manager backdates contracts into the current quarter. Which control BEST addresses this risk?
	\begin{enumerate}[label=\Alph*.,itemsep=1pt,topsep=3pt]
			\item Higher bonuses for quarter-end results.
			\item A reminder email about ethics.
			\item Faster commission calculation.
			\item Independent revenue-recognition review with cut-off testing of contract dates against delivery and approval records.
	\end{enumerate}
\end{frame}

\begin{frame}[t]{Q88 -- Explanation}
	\footnotesize
	\textbf{Correct: D.} Independent review and cut-off testing detect improper timing of revenue. \\[5pt]
\textbf{A wrong:} That increases incentive. \\[3pt]
\textbf{B wrong:} Reminders do not detect backdating. \\[3pt]
\textbf{C wrong:} Speed of commissions is unrelated. \\[3pt]
	\vfill {\tiny Set 2}
\end{frame}

% ===== Q89 =====
\begin{frame}[t]{Q89 \hfill {\scriptsize Fraud Risks}}
	\small
	\textbf{Question:} Which pattern is the MOST likely sign of overstated inventory?
	\begin{enumerate}[label=\Alph*.,itemsep=1pt,topsep=3pt]
			\item Rising turnover with rising sales.
			\item Inventory growing faster than sales, with falling turnover and no obsolescence write-downs.
			\item Stable inventory with stable sales.
			\item Falling inventory while sales rise.
	\end{enumerate}
\end{frame}

\begin{frame}[t]{Q89 -- Explanation}
	\footnotesize
	\textbf{Correct: B.} Disproportionate growth and falling turnover can signal overstated or obsolete inventory. \\[5pt]
\textbf{A wrong:} Consistent with healthy growth. \\[3pt]
\textbf{C wrong:} Consistent with stability. \\[3pt]
\textbf{D wrong:} Consistent with strong sales. \\[3pt]
	\vfill {\tiny Set 2}
\end{frame}

% ===== Q90 =====
\begin{frame}[t]{Q90 \hfill {\scriptsize Fraud Risks}}
	\small
	\textbf{Question:} Which is the strongest red flag for a bribery scheme?
	\begin{enumerate}[label=\Alph*.,itemsep=1pt,topsep=3pt]
			\item A competitive tender with documented evaluation.
			\item A modest documented training expense.
			\item A large success fee paid to an intermediary closely tied to a foreign official, with no documented services.
			\item A routine invoice from a long-established, unrelated supplier.
	\end{enumerate}
\end{frame}

\begin{frame}[t]{Q90 -- Explanation}
	\footnotesize
	\textbf{Correct: C.} Payments to connected intermediaries without clear services are classic bribery indicators. \\[5pt]
\textbf{A wrong:} Transparent and documented. \\[3pt]
\textbf{B wrong:} Normal and documented. \\[3pt]
\textbf{D wrong:} Routine and unconnected. \\[3pt]
	\vfill {\tiny Set 2}
\end{frame}

% ===== Q91 =====
\begin{frame}[t]{Q91 \hfill {\scriptsize Fraud Risks}}
	\small
	\textbf{Question:} Which analytic is MOST effective for detecting fictitious expense claims?
	\begin{enumerate}[label=\Alph*.,itemsep=1pt,topsep=3pt]
			\item Comparing total expenses with last year's total.
			\item Recomputing mileage rates.
			\item Verifying approvals only for the largest claims.
			\item Detecting duplicate or near-duplicate receipts and claims just under approval thresholds.
	\end{enumerate}
\end{frame}

\begin{frame}[t]{Q91 -- Explanation}
	\footnotesize
	\textbf{Correct: D.} Duplicates and threshold clustering are common signs of fabricated claims. \\[5pt]
\textbf{A wrong:} Aggregate comparisons miss individual manipulation. \\[3pt]
\textbf{B wrong:} Rates test accuracy, not authenticity. \\[3pt]
\textbf{C wrong:} Fraud often hides below thresholds. \\[3pt]
	\vfill {\tiny Set 2}
\end{frame}

% ===== Q92 =====
\begin{frame}[t]{Q92 \hfill {\scriptsize Fraud Risks}}
	\small
	\textbf{Question:} In a fraud investigation, when should the suspect typically be interviewed?
	\begin{enumerate}[label=\Alph*.,itemsep=1pt,topsep=3pt]
			\item First, to obtain a quick confession.
			\item Before securing records.
			\item Only if the board requests it.
			\item Near the end, after documents and witness evidence have been gathered.
	\end{enumerate}
\end{frame}

\begin{frame}[t]{Q92 -- Explanation}
	\footnotesize
	\textbf{Correct: D.} Interviewing the suspect last allows the investigator to be well prepared and reduces evidence-loss risk. \\[5pt]
\textbf{A wrong:} Early interviews tip off the suspect and weaken the case. \\[3pt]
\textbf{B wrong:} Records should be secured first. \\[3pt]
\textbf{C wrong:} The timing is an investigative decision, not a board request. \\[3pt]
	\vfill {\tiny Set 2}
\end{frame}

% ===== Q93 =====
\begin{frame}[t]{Q93 \hfill {\scriptsize Fraud Risks}}
	\small
	\textbf{Question:} Which is a detective control for cash handling fraud?
	\begin{enumerate}[label=\Alph*.,itemsep=1pt,topsep=3pt]
			\item A lock on the safe.
			\item Segregation of cash handling and recording.
			\item A dual-signature requirement.
			\item Surprise cash counts by an independent person.
	\end{enumerate}
\end{frame}

\begin{frame}[t]{Q93 -- Explanation}
	\footnotesize
	\textbf{Correct: D.} Surprise counts detect discrepancies after they occur. \\[5pt]
\textbf{A wrong:} Locks prevent access. \\[3pt]
\textbf{B wrong:} Segregation prevents concealment. \\[3pt]
\textbf{C wrong:} Dual signatures prevent unauthorized payments. \\[3pt]
	\vfill {\tiny Set 2}
\end{frame}

% ===== Q94 =====
\begin{frame}[t]{Q94 \hfill {\scriptsize Fraud Risks}}
	\small
	\textbf{Question:} Which fraud risk should the auditor assume exists in every organization?
	\begin{enumerate}[label=\Alph*.,itemsep=1pt,topsep=3pt]
			\item Theft by customers only.
			\item Competitor espionage only.
			\item Currency translation risk.
			\item Management override of controls.
	\end{enumerate}
\end{frame}

\begin{frame}[t]{Q94 -- Explanation}
	\footnotesize
	\textbf{Correct: D.} Management can override controls in any organization, so it must be considered in every assessment. \\[5pt]
\textbf{A wrong:} Not universal and not limited to customers. \\[3pt]
\textbf{B wrong:} Not a universal fraud risk. \\[3pt]
\textbf{C wrong:} That is a financial risk, not a fraud risk. \\[3pt]
	\vfill {\tiny Set 2}
\end{frame}

% ===== Q95 =====
\begin{frame}[t]{Q95 \hfill {\scriptsize Fraud Risks}}
	\small
	\textbf{Question:} Which pattern MOST suggests bid rigging?
	\begin{enumerate}[label=\Alph*.,itemsep=1pt,topsep=3pt]
			\item Many independent bidders with widely varying prices.
			\item The same few vendors rotate as lowest bidder and submissions contain identical errors.
			\item A single qualified vendor in a niche market, with documented justification.
			\item Winning bids that vary with the work requirements.
	\end{enumerate}
\end{frame}

\begin{frame}[t]{Q95 -- Explanation}
	\footnotesize
	\textbf{Correct: B.} Rotation and shared mistakes indicate collusion among bidders. \\[5pt]
\textbf{A wrong:} Independent prices suggest competition. \\[3pt]
\textbf{C wrong:} A documented sole-source can be legitimate. \\[3pt]
\textbf{D wrong:} Variation with requirements is normal. \\[3pt]
	\vfill {\tiny Set 2}
\end{frame}

% ===== Q96 =====
\begin{frame}[t]{Q96 \hfill {\scriptsize Fraud Risks}}
	\small
	\textbf{Question:} An employee who reported misconduct through the hotline is demoted shortly afterward. What should internal audit do?
	\begin{enumerate}[label=\Alph*.,itemsep=1pt,topsep=3pt]
			\item Treat it as unrelated to internal audit.
			\item Report the potential retaliation to the board or audit committee and recommend an independent review.
			\item Ask the manager who demoted the employee to investigate.
			\item Disclose the reporter's identity to affected managers.
	\end{enumerate}
\end{frame}

\begin{frame}[t]{Q96 -- Explanation}
	\footnotesize
	\textbf{Correct: B.} Retaliation undermines the reporting culture and should go to the governing body. \\[5pt]
\textbf{A wrong:} It affects the ethical culture and controls. \\[3pt]
\textbf{C wrong:} That manager has a conflict. \\[3pt]
\textbf{D wrong:} Revealing identity compounds the harm. \\[3pt]
	\vfill {\tiny Set 2}
\end{frame}

% ===== Q97 =====
\begin{frame}[t]{Q97 \hfill {\scriptsize Fraud Risks}}
	\small
	\textbf{Question:} Which control BEST prevents altered or forged company checks from being paid by the bank?
	\begin{enumerate}[label=\Alph*.,itemsep=1pt,topsep=3pt]
			\item Positive pay, where the bank matches presented checks to a list of issued checks.
			\item Monthly review of bank statements.
			\item A policy forbidding alterations.
			\item Annual employee ethics training.
	\end{enumerate}
\end{frame}

\begin{frame}[t]{Q97 -- Explanation}
	\footnotesize
	\textbf{Correct: A.} Positive pay stops unauthorized or altered checks before payment. \\[5pt]
\textbf{B wrong:} Review is detective and occurs after payment. \\[3pt]
\textbf{C wrong:} A policy alone does not stop clearing. \\[3pt]
\textbf{D wrong:} Training does not stop altered checks. \\[3pt]
	\vfill {\tiny Set 2}
\end{frame}

% ===== Q98 =====
\begin{frame}[t]{Q98 \hfill {\scriptsize Fraud Risks}}
	\small
	\textbf{Question:} A supposedly large vendor submits invoices with consecutive numbers over several months. What does this MOST likely indicate?
	\begin{enumerate}[label=\Alph*.,itemsep=1pt,topsep=3pt]
			\item Strong vendor controls.
			\item A system outage.
			\item The vendor may be fictitious or have few real customers.
			\item A volume discount.
	\end{enumerate}
\end{frame}

\begin{frame}[t]{Q98 -- Explanation}
	\footnotesize
	\textbf{Correct: C.} Sequential numbering from a purportedly large vendor suggests a shell or sham vendor. \\[5pt]
\textbf{A wrong:} Consecutive numbering is not evidence of control strength here. \\[3pt]
\textbf{B wrong:} Outages do not explain the pattern. \\[3pt]
\textbf{D wrong:} Discounts do not relate to numbering. \\[3pt]
	\vfill {\tiny Set 2}
\end{frame}

% ===== Q99 =====
\begin{frame}[t]{Q99 \hfill {\scriptsize Fraud Risks}}
	\small
	\textbf{Question:} A retail clerk fails to ring up a cash sale and keeps the money. Which scheme is this?
	\begin{enumerate}[label=\Alph*.,itemsep=1pt,topsep=3pt]
			\item Lapping.
			\item A billing scheme.
			\item Expense reimbursement fraud.
			\item Skimming.
	\end{enumerate}
\end{frame}

\begin{frame}[t]{Q99 -- Explanation}
	\footnotesize
	\textbf{Correct: D.} Skimming is taking cash before it is recorded. \\[5pt]
\textbf{A wrong:} Lapping conceals theft by applying later receipts to earlier accounts. \\[3pt]
\textbf{B wrong:} Billing schemes involve false invoices. \\[3pt]
\textbf{C wrong:} No expense claim is involved. \\[3pt]
	\vfill {\tiny Set 2}
\end{frame}

% ===== Q100 =====
\begin{frame}[t]{Q100 \hfill {\scriptsize Fraud Risks}}
	\small
	\textbf{Question:} Which factor MOST strongly deters employee fraud?
	\begin{enumerate}[label=\Alph*.,itemsep=1pt,topsep=3pt]
			\item A strong perception that misconduct will be detected and result in consequences.
			\item A longer code of conduct.
			\item A higher salary for all staff.
			\item An annual fraud video with no testing.
	\end{enumerate}
\end{frame}

\begin{frame}[t]{Q100 -- Explanation}
	\footnotesize
	\textbf{Correct: A.} Perceived likelihood of detection and consequences is a powerful deterrent. \\[5pt]
\textbf{B wrong:} Length does not deter without enforcement. \\[3pt]
\textbf{C wrong:} Pay does not remove opportunity or rationalization. \\[3pt]
\textbf{D wrong:} Passive awareness has limited effect. \\[3pt]
	\vfill {\tiny Set 2}
\end{frame}

\begin{frame}{Answer key}
\small
1:A \quad 2:A \quad 3:B \quad 4:A \quad 5:D \quad 6:C \quad 7:B \quad 8:C \quad 9:B \quad 10:A \quad 11:D \quad 12:A \quad 13:C \quad 14:D \quad 15:D \quad 16:D \quad 17:B \quad 18:A \quad 19:B \quad 20:A \quad 21:C \quad 22:B \quad 23:A \quad 24:B \quad 25:B \\[6pt]
26:C \quad 27:A \quad 28:D \quad 29:C \quad 30:B \quad 31:B \quad 32:B \quad 33:A \quad 34:D \quad 35:C \quad 36:C \quad 37:D \quad 38:C \quad 39:A \quad 40:D \quad 41:B \quad 42:D \quad 43:A \quad 44:C \quad 45:C \quad 46:B \quad 47:A \quad 48:C \quad 49:B \quad 50:B \\[6pt]
51:B \quad 52:A \quad 53:D \quad 54:C \quad 55:C \quad 56:A \quad 57:B \quad 58:A \quad 59:D \quad 60:C \quad 61:A \quad 62:B \quad 63:A \quad 64:D \quad 65:D \quad 66:C \quad 67:D \quad 68:C \quad 69:A \quad 70:C \quad 71:C \quad 72:D \quad 73:A \quad 74:C \quad 75:D \\[6pt]
76:A \quad 77:D \quad 78:C \quad 79:C \quad 80:B \quad 81:B \quad 82:D \quad 83:C \quad 84:A \quad 85:B \quad 86:A \quad 87:B \quad 88:D \quad 89:B \quad 90:C \quad 91:D \quad 92:D \quad 93:D \quad 94:D \quad 95:B \quad 96:B \quad 97:A \quad 98:C \quad 99:D \quad 100:A
\end{frame}
\end{document}
